% 本文件由 LaTeX 排版 Agent 根据有效 translation.json 自动生成。
% author=Augustine of Hippo; work_id=1507; title=论婚姻与贪欲

\chapter{论婚姻与贪欲}

% structure: p0001 source=h1 target=omitted reason=page-title-already-rendered-by-parent

两卷，致\Person{瓦莱里乌斯伯爵}\par

\Person{奥古斯丁}，\Place{希波}主教；写于419及420年间\par

\SourceRecord{Translated by Peter Holmes and Robert Ernest Wallis, and revised by Benjamin B. Warfield.；Nicene and Post-Nicene Fathers, First Series；Vol. 5.；Edited by Philip Schaff.；Buffalo, NY: Christian Literature Publishing Co.,；1887.；<http://www.newadvent.org/fathers/1507.htm>.}

% structure: p0001 source=h1 target=section reason=page-title

\section{论婚姻与贪欲（附函）}

\EditorialNote{致瓦莱里乌斯伯爵的一封信，论及奥古斯丁向他寄送他称为第一本书的《论婚姻与贪欲》。}\par

给在基督的爱中杰出的、应受尊敬的主上、最亲爱的儿子瓦莱里乌斯，奥古斯丁在主内问候。\par

1. 我正在为许久未收到殿下对我所写多封信件的回复而感到懊恼时，突然收到了您恩赐的三封信——一封由我的同为主教的文德米亚利斯带来，并非专指给我；另两封不久后由我的弟兄、司铎菲尔姆斯带来。这位圣人，您可能从他本人那里得知，他以最亲密的爱与我相连，曾与我多次谈论您的卓越，并向我提供了无可置疑的证据，表明他完全了解您的品格\Quote{在基督的怀抱内；}\BibleRef{斐 1:8} 通过这些方式，他不仅看到了前述主教和他本人所持的信件，也看到了那些我们未收到而表示失望的信件。他对您的了解令我们格外欣慰，因为他让我明白，您出于您的权力无法做到的事——即使我恳切请求回信，您也不会成为自己赞美的宣扬者，这与圣经的许可相反\BibleRef{箴 27:2}。但我自己也应犹豫是否以这样的笔调给您写信，以免招致谄媚之嫌，我杰出且应受尊敬的、在基督的爱中最亲爱的儿子主上。\par

2. 至于您在基督内的赞美，或不如说基督在您内的赞美，您看我从他那里听闻这些时是多么欣喜和快乐！他因对我的忠诚不会欺骗我，又因与您的友谊不会对它们无知。但其他见证，虽在数量和确定性上较差，仍从各方传到我耳中，向我保证您的信德是多么纯正和公教，对未来的望德是多么虔诚，对天主和弟兄的爱德多么伟大，身处最高荣耀中仍多么谦卑，因为您不倚靠不定的财富，而倚靠永生的天主，并在善工上富足\BibleRef{弟前 6:17}；您的家是圣徒的安息和安慰，是恶人的恐惧；您多么留意不让人（无论是祂的旧敌还是新近的敌人）用基督的名作为诡计的外衣，为基督的肢体设下圈套；同时，尽管您对这些敌人的错误毫不留情，却何其谨慎地确保他们的得救。凡此种种，正如我已说的，我们也常从其他人口中听到；但此刻，通过上述弟兄，我们获得了更全面、更可靠的认识。\par

3. 然而，关于婚姻贞洁，为了我们能为此称赞和爱您，我们怎么可能只听从任何除了您某位心腹朋友之外的人的讲述？这位朋友对您不是肤浅的了解，而是深知您的内心生活。因此，关于天主赐予您的这一美好恩赐，我乐于更坦率、更详尽地与您交谈。我确信，即使我寄给您一篇冗长的论文，其阅读只会确保我们之间更长的对话，我也不会成为您的负担。因为我发现，在您众多繁重的忧虑中，您轻松愉快地阅读，并且，每次我们的一些小作品（即使寄给他人）偶然落入您手中，您都从中获得极大乐趣。因此，凡是寄给\Emph{您自己}的，我能像当面与您交谈那样对您说话，您将屈尊以更大的关注留意它，并以更高的喜悦接受它。那么，在阅读完这封信后，请转向我随信附上的书。该书在其开篇将以更合适的方式向您尊驾表明它写作的原因，以及为何特别呈交您审阅。\par

\SourceRecord{Translated by Peter Holmes and Robert Ernest Wallis, and revised by Benjamin B. Warfield.；Nicene and Post-Nicene Fathers, First Series；Vol. 5.；Edited by Philip Schaff.；Buffalo, NY: Christian Literature Publishing Co.,；1887.；<http://www.newadvent.org/fathers/15070.htm>.}

% structure: p0001 source=h1 target=section reason=page-title

\section{论婚姻与贪欲（卷一）}

\EditorialNote{本书阐明婚姻特有的自然祝福。他指出，在这些祝福中，不可将肉体的贪欲计算在内，因为它完全是恶的，并非来自婚姻的本质，而是源于原罪的一种意外。不过，婚姻仍可正当地利用此恶来生育子女。但由于这种贪欲的结果，即使在天主子女的合法婚姻中，所生的孩子也不是天主的子女，而是世界的子女，被罪的锁链捆绑，尽管他们的父母已因恩宠而获解放；并且被魔鬼俘虏，若他们不被基督的同一恩宠同样拯救。他解释贪欲如何在已受洗者身上存留于行为而非罪责。他教导，藉着圣洗圣事的神圣性，不仅原罪，而且人的所有其他罪都被除去。最后他引用盎博罗削的权威来表明贪欲之恶必须与婚姻之善区分开来。}\par

% structure: p0003 source=h2 target=subsection reason=relative-heading-rank; base=h2

\subsection{第一章——论本书的主旨}

我们亲爱的儿子\Person{瓦莱里乌斯}，我们的新异端者坚持认为在肉身中出生的婴儿不需要基督的药来医治罪，他们出于对我们的极度仇恨，不断断言我们谴责婚姻以及天主通过男女生育人类的神圣程序，因为我们断言从这种结合出生的人沾染了原罪，正如宗徒所说：\Quote{藉着一个人，罪恶进入了世界，死亡藉着罪恶而来；于是死亡临到了众人，因为众人都犯了罪}；并且我们不否认，无论他们是什么样的父母所生，他们仍处于魔鬼的统治之下，除非他们在基督里重生，并藉着他的恩宠从黑暗的权势中被救出，迁入他的国度\BibleRef{哥 1:15}——他愿意不从两性的结合中出生。因为，我们肯定这个包含在最古老且不变的公教信仰准则中的教义，而这些提出新奇且乖僻教义的人断言婴儿在重生之洗\BibleRef{弟 3:5}中没有需要被洗去的罪，他们因不信或无知而诽谤我们，好像我们谴责婚姻，并且断言天主的作为——由婚姻而生的人——是魔鬼的工作。他们也不反思，婚姻的善不能因从它而来的原恶而被指责，正如奸淫和私通的恶不能因从它们而生的自然善而被原谅。因为，罪是魔鬼的工作，无论婴儿从哪里染上；而人是天主的工作，无论从哪里出生。因此，本书的目的，在主恩准我们帮助的范围内，是要区分肉体的贪欲之恶（由此出生的人染上原罪）与婚姻之善。因为，如果没有人类先前犯罪，这种可耻的贪欲就不会存在，它被无耻的人无耻地赞美；而婚姻，如果没有人犯罪，它仍然会存在，因为在属于那生命的身体中，生育子女将不会伴随这存在于\Quote{这个死亡的身体}\BibleRef{罗 7:24}中无法与生育过程分离的疾病。\par

% structure: p0005 source=h2 target=subsection reason=relative-heading-rank; base=h2

\subsection{第二章 [II.]——为何此书致信瓦莱里乌斯}

现在有三个非常特殊的原因，我将简要说明，为什么我特别想给您写这封信。一是，由于基督的恩赐，您是婚姻贞洁的严格持守者。二是，由于您的巨大关心和勤奋，您有效地抵挡了我们当前讨论中对抗的那些亵渎新奇事物。三是，因为我了解到他们写成的一些东西曾落入您手中；虽然以您坚强的信德，您可以鄙视这种企图，但通过辩护来帮助我们的信德仍是一件好事。因为宗徒\Person{伯多禄}教导我们要\Quote{时常准备好答复每一个问我们心中所怀信德和望德理由的人}；宗徒\Person{保禄}说：\Quote{你们的言谈要常带着恩宠，用盐调和，好知道怎样答复每一个人。}\BibleRef{哥 4:6}这些是主要激励我在主所允许的范围内，在这卷书中与您进行这样的对话的动机。我确实从未喜欢过任何一位像您这样因其职位高度而众人皆知的杰出人士，在没有他请求的情况下，将我的任何谦卑劳作的阅读强加给他——尤其当他没有享受尊贵的退休生活，而是仍在从事士兵职业的公共职责时；这在我看来更像莽撞而非尊重。如果，在我现在提到的理由下行动时，我招致了这种指责，我恳请您原谅并友好地看待以下论证。\par

% structure: p0007 source=h2 target=subsection reason=relative-heading-rank; base=h2

\subsection{第三章 [III.]——婚姻贞洁是天主的恩赐}

婚姻中的贞洁是天主的恩赐，这由最蒙福的\Person{保禄}在谈到这主题时指出：\Quote{我愿众人都像我一样，但每人各自从天主那里领受自己的恩赐，有人这样，有人那样。}\BibleRef{格前 7:7}请注意，他告诉我们这恩赐来自天主；尽管他将其列于他愿众人像他一样的该节制之下，他仍将其描述为天主的恩赐。由此我们明白，当这些诫命赐给我们以便我们遵行时，没有别的意思，只是说我们内也应有自己的意志去接受并拥有它们。因此，当这些被显示为天主的恩赐时，意思是如果还没有拥有，就必须向他寻求；如果已经拥有，就必须为拥有而感谢他；再者，我们自己的意志在寻求、获得和持守这些恩赐方面几乎没有作用，除非它们被天主的恩宠所助佑。\par

% structure: p0009 source=h2 target=subsection reason=relative-heading-rank; base=h2

\subsection{第四章——关于不信者贞洁的难题。只有信徒才是真正贞洁的人。}

那么，当我们在某些不信者身上也发现婚姻贞洁时，我们该说什么呢？是否必须说他们犯罪，因为他们滥用了天主的恩赐，没有将这恩赐归还原主，即赐予他们的那位的敬拜呢？或者，当行使这些恩赐的人并非信徒时，或许这些禀赋根本不应被视为天主的恩赐？这是按照宗徒的说法，他说：\BibleQuote{凡不出于信德的都是罪？}\BibleRef{罗 14:23}但谁敢说天主的恩赐是罪呢？因为灵魂和身体，以及一切植入灵魂和身体中的本性禀赋，即使在罪人身上，仍然是天主的恩赐；因为是天主造了他们，而不是他们自己造的。当说\BibleQuote{凡不出于信德的都是罪，}仅指人自己所做的事。因此，当人没有信德而做那些似乎属于婚姻贞洁的事时，他们这样做或是为了取悦人（无论是自己或他人），或是为了避免陷入人性在其所邪欲之事中容易遭遇的烦恼，或是为了侍奉魔鬼。罪并非真正被放弃，而是某些罪被另一些罪所压倒。所以，愿天主不许人被称为真贞洁，如果他对妻子保持婚姻忠诚并非出于对真天主的虔诚。\par

% structure: p0011 source=h2 target=subsection reason=relative-heading-rank; base=h2

\subsection{第五章 [IV.]——婚姻的自然之善。一切社会自然拒绝欺诈的伴侣。什么是真正的婚姻贞洁？除非奉献给真信仰，否则没有真正的童贞和贞洁。}

那么，男女为生育而结合是婚姻的自然之善。但谁若兽性地使用这一善，使其意图在于满足情欲而非渴望子嗣，便是滥用此善。然而，在许多无理性的动物中，例如大多数鸟类，既保存着一种成对结合，又有筑巢技能的社会合作；它们相互分工抚卵期，交替喂养幼雏，这使得它们在交配时显得更像是致力于确保物种延续而非满足情欲。二者之中，一个是人借禽兽的形像；另一个是禽兽借人的形像。然而，关于我所归于婚姻本性的，即男女作为生育的伴侣而结合，因此不互相欺骗（因为每一种结合状态都天然憎恶欺诈的伴侣），尽管不信的人也拥有这明显的本性祝福，但由于他们不信德地使用它，反而将其转为罪恶。因此，同样地，信徒的婚姻将那肉欲（由此\BibleQuote{肉情相反圣神}）转为正义之用。\BibleRef{迦 5:17}因为他们怀着坚定的意愿生育子女，使之得以重生——即他们生为\Quote{今世之子}的子女可以重生而成为\Quote{天主的子女}。因此，凡不以此意向、意愿、目的——即将子女从第一人的肢体转变为基督的肢体——来生育子女的父母，反而像不信的父母一样以不信的子女夸耀——无论他们在同居中多么谨慎，刻意将之限于生育子女——实际上自身并没有婚姻贞洁。因为贞洁是德行，憎恶不洁作为相反的恶习，而所有德行（甚至那些通过身体运作的）都居于灵魂之中，那么当灵魂自身在对真天主行邪淫时，身体怎能真正被称为贞洁呢？这种邪淫正是圣咏作者所谴责的，他说：\Quote{看，那远离祢的必灭亡；凡离开祢行淫的，祢都消灭了。}因此，没有真正的贞洁，无论是婚姻的、寡居的、还是童女的，除非是将自己奉献给真信仰。因为虽然献身的童贞被正当地置于婚姻之上，但哪个头脑清醒的基督徒不会宁可选择那些甚至结过多次婚的公教基督徒妇女，而不是那些灶神贞女，甚或异端童女呢？信德的效果何其大！关于它，宗徒说：\BibleQuote{凡不出于信德的都是罪；}\BibleRef{罗 14:23}又在希伯来书中写道：\BibleQuote{没有信德，不可能中悦天主。}\BibleRef{希 11:6}\par

% structure: p0013 source=h2 target=subsection reason=relative-heading-rank; base=h2

\subsection{第六章 [V.]——谴责情欲并非谴责婚姻；人身体中的羞耻从何而来。亚当和厄娃受造时并非盲者；他们\Quote{眼睛开了}的意义。}

那么，既然问题的真相如此，那些认为当肉欲受到谴责时婚姻也被定罪的人无疑错了；好像贪欲的病症是婚姻的结果而不是罪的结果。难道天主用\BibleQuote{你们要生育繁殖}\BibleRef{创 1:28}的话祝福的第一对夫妇，他们不是赤身露体却不羞耻吗？那么，为何在犯罪后羞耻从他们的肢体中产生呢？除非是因为一种不雅的冲动从它们中产生，如果人没有犯罪，婚姻中肯定绝不会存在这种冲动。或者，正如某些人所认为的（他们不太留意自己所读的），人类起初像狗一样被造为瞎眼的，而更荒谬的是，他们不是像狗那样通过成长，而是通过犯罪得到视力？我们绝不可有此想法。但他们从读经中得出这个观点：\BibleQuote{她就摘了那树的果子，吃了，也给了与她在一起的丈夫，他也吃了；二人的眼就开了，才知道自己是赤身露体的。}\BibleRef{创 3:6}这解释了那些无知之人的看法，即第一个男人和女人的眼睛先前是闭着的，因为圣经证实它们当时被打开了。那么，\Person{撒辣}的婢女\Person{哈加尔}的眼睛先前是闭着的吗？当她与口渴哭泣的孩子一起时，她开了眼\BibleRef{创 21:17}看见了水井？或者，那两个门徒在主的复活后，是否闭着眼睛与他同行？因为福音作者论到他们说\BibleQuote{在擘饼时他们的眼睛开了，就认出他来}\BibleRef{路 24:31}？因此，论到第一个男人和女人所写的\BibleQuote{二人的眼就开了}\BibleRef{创 3:7}，我们应当理解为：他们开始注意觉察并认出他们身体所遭的新状态。既然他们的眼睛开了，他们的身体就向他们显得赤露，他们就知道了。如果这不是含义，那么，当日间走兽和空中飞鸟被带到\Person{亚当}面前时\BibleRef{创 2:19}，如果他的眼睛是闭着的，他怎能给它们命名呢？他不可能不区分它们就做到这一点；而他不可能不看见它们就区分它们。再者，当他这样说时，他怎能如此清晰地看见那女人：\BibleQuote{这才是我的骨中之骨，肉中之肉}\BibleRef{创 2:23}？如果真有人好辩到坚持认为\Person{亚当}可能不是靠视觉而是靠触觉获得对这些对象的辨识，那么他将如何解释那段经文，其中告诉我们女人\BibleQuote{看那棵树，}她正要摘取禁果，\BibleQuote{实在好看}\BibleRef{创 3:6}？不；\BibleQuote{他们二人都赤身露体，并不羞耻}\BibleRef{创 2:25}，不是因为他们没有视力，而是因为他们没有察觉到自己的肢体中有任何可羞耻的理由，这些肢体一直为他们所见。因为并没有说：他们二人都赤身露体，\Emph{而不知道}；而是说\BibleQuote{他们并不羞耻}。因为，确实，先前没有发生任何不合法的事，因此也没有产生任何能使他们羞耻的事。\par

% structure: p0015 source=h2 target=subsection reason=relative-heading-rank; base=h2

\subsection{第七章 [第六篇]——人的悖逆在其自身肉体的反叛中受到正义的报应；因身体悖逆的肢体而脸红羞愧。}

当第一个人违犯了天主的法律，他开始在肢体中有另一条法律，与他的理智法律相敌，他感到自己悖逆的邪恶，当他体验到自己肉体悖逆中最正义的报应反弹到自己身上时。那么，这就是蛇在诱惑中所应许的\BibleQuote{他眼睛的开启}\BibleRef{创 3:5}——事实上，是对他宁愿不知道的某件事的知识。那时，人的确在自己内察觉他所做了的事；那时他区分善恶——不是通过避免恶，而是通过忍受恶。因为确实不应当，他的仆人，即他的身体，应服从他，而他却没有服从他自己的主。那么，这个事实何其意味深长：眼睛、嘴唇、舌头、手、脚、背脊的弯曲、颈项、腰肋，全都置于我们的能力之内——当我们的身体无障碍且健康时，可用来进行适合它们的动作；但当涉及人类生育子女的伟大功能时，明确为此目的而创造的肢体却不服从意志的指示，而必须等待肉欲来使这些肢体运动，好像肉欲对它们有法定权利似的，而且有时当意念愿意时它拒绝行动，而常常它又违反意志而行动！这难道不使人类自由意志因藐视天主、自己的统帅，从而丧失对自己的肢体的一切正当支配而羞愧脸红吗？那么，何处能找到一个更合适的证明，表明人性因其悖逆而受到的正义败坏，比那些部分——自然本身通过它们得以依次延续——的悖逆更合适呢？因为身体那些部分被特别恰当地称为\Emph{自然的}。那么，这就是为何第一对人类夫妇，在肉体中体验到那因悖逆而不雅的冲动，并感到自己赤身露体的羞耻后，用无花果叶遮盖了这些犯罪的肢体\BibleRef{创 3:7}；为的是，至少，通过羞耻的犯罪者的意志，一个面纱可以抛在那些不按希望者的意志而被推动的东西上；并且，既然羞耻来自于那下流地取悦的东西，体面可以通过隐藏而获得。\par

% structure: p0017 source=h2 target=subsection reason=relative-heading-rank; base=h2

\subsection{第八章 [第七篇]——肉欲的邪恶并不除去婚姻的美善。}

既然婚姻之善不会因这恶的加入而丧失，有些轻率之人便以为这并非加添的恶，而是属于原初之善。然而，不仅精微的理性，就连最普通的自然判断也能区分，这在第一个男人和女人身上显而易见，如今在已婚者身上依然成立。他们后来在生育中所成就的——那是婚姻之善；但他们最初因羞耻而遮掩的——那是\OriginalTerm{贪欲}{concupiscentia}之恶，这恶处处躲避视线，在羞耻中寻求隐秘。因此，婚姻既从这恶中生出一些善，便有可夸之处；但既然善不能脱离恶而实现，它便有理由感到羞耻。这可以藉由瘸子的例子来说明。假设他跛行着达到某个善的目标，那么，一方面，成就本身并不因人的跛行之恶而为恶；另一方面，跛行也不因成就之善而为善。同理，我们不应因淫欲之恶而谴责婚姻，也不应因婚姻之善而赞美淫欲。\par

% structure: p0019 source=h2 target=subsection reason=relative-heading-rank; base=h2

\subsection{第九章 [VIII.]——这贪欲之病在婚姻中不应出于意愿，而应出于必要；信徒在婚姻中使用时应如何意愿；何谓使用而非屈服于贪欲之恶；旧约的圣祖们从前如何使用妻子。}

宗徒所指的正是这贪欲之病，当他向已婚信徒说话时说：\BibleQuote{这就是天主的旨意，你们的圣化，即你们应远避邪淫：让每个人都知道，应以圣化和尊敬持守自己的器皿，而不在欲望之病中，如同不认识天主的外邦人。} \BibleRef{得前4:3} 因此，已婚信徒不仅不可使用他人的器皿，即那些贪恋他人妻子的人所行的；而且必须知道，连自己的器皿也不可在肉体的贪欲之病中持守。这劝谕不应理解为宗徒禁止夫妇——即合法而可敬的——同居；而是说这种同居（若非因先前的罪使人完全的自由选择变得如此无能，以致有了这致命的附赘，它本不会有不健康的情欲伴随）不应出于意愿，而应出于必要，然而若没有它，便不可能实现生育子女的意愿本身。信徒婚姻中的这一意愿并非由让子女在此世度过一生的目的所决定，而是为使他们在基督内重生，并永远留在祂内。如果这一结果得以实现，圆满幸福的赏报将从婚姻而来；但若不能实现，夫妇之间仍会有善意带来的平安。凡以此心怀持守自己器皿（即妻子）的人，肯定不会像不认识天主的外邦人那样在\BibleQuote{欲望之病}中持守她，而是像仰望天主的信徒那样，在圣化和尊敬中持守。当人约束并抑制贪欲之恶的狂怒，它活动在无节制、不当的运动中时，他是在使用这恶而不被它征服；除非为了子嗣，他绝不放松对它的控制，然后引导并应用它于生育将在灵性上重生的子女，而不是使灵性屈服于肉体，陷入污秽的奴役。亚巴郎之前和之后的古代圣祖们，天主曾为他们作证说\BibleQuote{他们中悦了祂}，他们就是这样使用自己的妻子，任何基督徒都不应怀疑，因为其中某些人曾被允许有多个妻子，其原因是为了多生子嗣，而非追求不同满足的欲望。\par

% structure: p0021 source=h2 target=subsection reason=relative-heading-rank; base=h2

\subsection{第十章 [IX.]——为何有时允许一个男人有多个妻子，却从未允许一个女人有超过一个丈夫；本性在它的\OriginalTerm{主权}{dominationes}中偏好单一。}

如今，如果对我们祖先的天主，也就是我们的天主，并不厌恶如此多的妻子，目的是使情欲有更充分的放纵余地；那么，按此假设，圣善的妇女也应该各自服侍多个丈夫。但若任何妇女这样做，除了可耻的情欲之外，还有什么情感能促使她拥有多个丈夫？因为这样的自由并不能使她生育更多子女。然而，婚姻的良善目的，由一夫一妻比一夫多妻更能实现，这已由天主亲自缔结的最初婚姻结合充分显明。祂有意使婚姻从此开始，并为它们提供一个更可敬的先例。然而，随着人类的进展，某些良善的男子被赐予了多位良善的妻子——每一位都有许多；由此看来，节制似乎为了尊严而寻求一方的单一，而为了多产，本性却允许另一方的众多。因为按本性之理，一人统治多人比多人统治一人更可行。同样不可怀疑的是，男人统治女人比女人统治男人更合乎本性的秩序。宗徒正是本着这一原则说：\BibleQuote{男人是女人的头；}\BibleRef{格前 11:3} 以及\BibleQuote{妻子们，要服从自己的丈夫。}\BibleRef{哥 3:18} 同样，宗徒伯多禄也写道：\BibleQuote{就如撒辣服从了亚巴郎，称他为主。}\BibleRef{伯前 3:6} 然而，尽管事实上，本性在其统治中喜爱单一，但我们可能在人类中较低下的部分更容易看到多数；尽管如此，除非为了从他生出更多的子女，一个人拥有多个妻子从来不是合法的。因此，如果一个女人与多个男人同居，既然她不能因此增加后代，而只是更频繁地满足情欲，她就不能是妻子，而只能是娼妓。\par

% structure: p0023 source=h2 target=subsection reason=relative-heading-rank; base=h2

\subsection{第十一章 [X.]——婚姻的圣事；婚姻不可解散；世间关于离婚的法律与福音的不同。}

婚姻当然不只是多产（其果实在于后代），也不只是贞洁（其纽带是忠信），而且还有某种圣事性的纽带，它为信徒的婚姻所推荐。因此宗徒命令说：\BibleQuote{丈夫们，要爱你们的妻子，如同基督也爱了教会。}\BibleRef{弗 5:25} 这纽带的实质无疑是：在婚姻中结合的男女应当终身不可分离；除非因奸淫的缘故，一方不得与另一方分离。\BibleRef{玛 5:32} 因为这保存在基督与教会的情形中；这样，如同活人与活人，没有离婚，永远没有分离。在我们天主的城、祂的圣山上——就是说，在基督的教会——所有已婚的信徒（他们无疑是基督的肢体）如此完全遵守这一纽带，以至于虽然男女结婚是为了生育子女，但绝不允许一个人为了另娶而休弃甚至不生育的妻子。凡这样做的人，按福音的法律被视为犯奸淫罪；虽然按世间的规则并非如此，因为世间的规则允许双方离婚，甚至无需指控有罪，并可缔结新的婚约——主告诉我们，连圣梅瑟也因他们心硬而允许了以色列人民这样做。\BibleRef{玛 19:8} 同样，女人若另嫁他人，也受同样的定罪。确实，婚姻的权利在缔结婚姻者之间如此持久，只要双方都活着，甚至彼此分离的人仍被视为夫妻，而不与另结新欢的人相同。因为若先前的婚姻关系不继续存在，他们就不会因这新结合而犯奸淫。如果丈夫死了，与他缔结了真婚姻，那么现在便可以缔结真婚姻，而这在以前会是奸淫。因此，夫妻之间，只要活着，婚姻纽带就有永久的约束，既不能因分离而取消，也不能因与他人结合而取消。但这种持久性在这种情况下只对罪的伤害有用，而不是盟约的纽带。同样，背教者的灵魂，它仿佛放弃了与基督的婚姻结合，即使它抛弃了信德，也不会失去它在重生之洗中所领受的信德的圣事。如果他回归，这圣事无疑会归还给他，虽然他离开基督时失去了它。然而，他在背教之后仍保留着圣事，这加重了他的惩罚，而不是获得赏报。\par

% structure: p0025 source=h2 target=subsection reason=relative-heading-rank; base=h2

\subsection{第十二章 [XI.]——婚姻不取消互守贞洁的誓愿；玛利亚和若瑟之间确有真婚姻；若瑟如何是基督的父亲。}

但天主绝不会允许将那些因彼此同意而决定永远戒绝肉体私欲偏情之人的婚姻纽带视为破裂。不，它只会更加牢固，因为他们交换了誓约，这誓约须以特殊的情感与和谐来持守——并非藉肉体的纵欲纽带，而是藉灵魂的自愿情感。因为天使并非虚假地对若瑟说：\BibleQuote{不要怕娶你的妻子玛利亚。}\BibleRef{玛 1:20}她之所以被称为他的妻子，是因她最初的婚约，尽管他从未与她有肉体的认识，也永不会有。妻子这个称谓既未被摧毁，也非虚假，尽管从未有、且无意有肉体的结合。这位童贞妻子对她丈夫而言，反而因其无男人而怀孕，成为更神圣、更奇妙的喜乐，虽在所生之子方面有差异，但在他们所怀的信德上并无差异。因这婚姻的信实，他们两人都当得起被称为基督的\BibleQuote{父母}\BibleRef{路 2:41}（不仅她作为祂的母亲，他也作为祂的父亲，因他是她的丈夫），二者在心志与意向上是如此，虽非在肉身上。然而，前者（若瑟）仅在意志上是祂的父亲，后者（玛利亚）也在肉体上是祂的母亲，但他们二者都只是祂谦卑的父母，并非祂崇高的父母；是祂软弱的父母，并非祂神性的父母。因为福音并不说谎，其中写道：\Quote{祂的母亲和父亲都惊奇那些关于祂所说的话；}又在另一处说：\BibleQuote{祂的父母每年上耶路撒冷去；}\BibleRef{路 2:41}稍后又写道：\BibleQuote{祂的母亲对祂说：\InnerQuote{孩子，你为什么这样对待我们？看，你的父亲和我，一直忧苦地找你。}}\BibleRef{路 2:48}然而，为表明祂另有一位超越他们的父亲，那位无母而生祂的父，祂回答他们说：\BibleQuote{你们为什么找我？你们不知道我必须在我父那里吗？}\BibleRef{路 2:49}再者，免得有人认为祂因刚才所说的话而否认他们是祂的父母，福音圣史立刻加上：\BibleQuote{他们对祂所说的话不明白；祂就同他们下去，来到纳匝肋，属他们管辖。}\BibleRef{路 2:50}属谁管辖？岂非祂的父母？谁在属他们管辖？岂非耶稣基督，\BibleQuote{祂虽具有天主的形体，却没有把与天主同等视为应紧紧把持的，}\BibleRef{斐 2:6}？祂为何要服从那些远低于天主形体的人？除非是因为\BibleQuote{祂空虚自己，取了奴仆的形体}\BibleRef{斐 2:7}——即祂父母所处的形体。既然她生祂时没有丈夫的配合，他们二人不可能同为那奴仆形体的父母，除非他们虽无肉体的结合，却已因婚姻而结合。因此，世系系列（尽管基督的双亲在谱系中一同被提及）\BibleRef{玛 1:16}必须，且实际上已经，延续到若瑟的名字为止，以免在这婚姻中愧对男性——而且是较强壮的性别——同时不损害真理，因为若瑟与玛利亚一样，都出自达味的后裔，\BibleRef{路 1:27}预许说基督当从此而来。\par

% structure: p0027 source=h2 target=subsection reason=relative-heading-rank; base=h2

\subsection{第十三章——在玛利亚和若瑟的婚姻中，婚姻的全部福祉俱备；凡生于姘居者皆为有罪的肉躯。}

因此，婚姻制度的全部好处已在基督的这双亲身上实现：有子嗣，有忠信，有约束。作为子嗣，我们认出主耶稣本身；忠信在于没有奸淫；约束在于没有休妻。【十二】只是没有婚配同居；因为那要成为无罪的、且被派遣\BibleQuote{成为罪过的肉躯的\Emph{相似体}}\BibleRef{罗 8:3}的祂，不可能在本身为有罪的肉躯中受生，而不经历那来自罪、可耻的肉体私欲偏情；祂不愿带着这私欲偏情而生，为教导我们：凡生于性交者实际上都是有罪的肉躯，因为惟独那非由性交所生者才不是有罪的肉躯。尽管如此，婚配的同居本身并非罪恶，当它以生育子女为意图而进行时；因为心灵之善意引领随之而来的肉体快感，而非追随其主导；人的抉择并不因罪轭的压迫而分心，因为罪的打击适当地被引向生育的目的。这打击具有某种挑逗性的活动，在奸淫、私通、放荡和不洁的臭恶纵欲中称王为霸；而在婚姻状态不可或缺的职责中，它却表现出奴仆的顺从。在前者，它被谴责为暴虐之主厚颜无耻的放肆；在后者，它作为如此温顺的仆从的忠实服务而得到适度的赞扬。因此，这私欲偏情本身并非婚姻制度之善；而是有罪之人中的污秽、生育父母中的必要、纵欲的火焰、婚配欢愉的羞耻。那么，为什么当夫妻双方同意停止同居时，他们不能继续保持夫妻关系呢？既然若瑟和玛利亚从未开始同居，却始终是夫妻。\par

% structure: p0029 source=h2 target=subsection reason=relative-heading-rank; base=h2

\subsection{第十四章——基督以前是婚配之时；自基督以来是守节之时。}

如今，这种子女的繁衍在古代圣人间曾是一项最本分的义务，目的在于为天主生育并保存一个民族，在这民族中，基督来临的预言必须优先于一切；如今，它已不再具有同样的必要性。因为从所有民族中，无论他们从何处获得自然出生，都有丰富的后代可获得属神的再生。这样我们就能明白圣经所说的：\BibleQuote{拥抱有时，戒避拥抱有时}\newline{}\newline{}\BibleRef{训 3:5}，这句话的各个分句应分别属于基督之前和基督之后的时期。前者是拥抱的时期，后者是戒避拥抱的时期。\par

% structure: p0031 source=h2 target=subsection reason=relative-heading-rank; base=h2

\subsection{第十五章——宗徒关于此事的教导。}

因此，宗徒似乎也看到这段经文，他声明说：\BibleQuote{但是，弟兄们，我给你们说：时间是短促的；从此以后，有妻子的要像没有一样；哭泣的要像不哭泣的；欢乐的要像不欢乐的；购买的要像一无所得的；享用这世界的要像不享用的，因为这世界的局面正在逝去。但我愿意你们无所挂虑。}\newline{}\newline{}\BibleRef{格前 7:29} 整个这段经文（为了就这个主题表达我的看法，我以对宗徒言语的简短诠释形式出现），我认为必须理解如下：\BibleQuote{弟兄们，我给你们说：时间是短促的。}天主的子民不再藉着肉体的生育来繁衍；而是从此以后，要藉着属神的再生聚集起来。\BibleQuote{所以，从此以后，有妻子的}不应受肉欲的支配；\BibleQuote{哭泣的}面对当前邪恶的悲伤，应因未来福佑的希望而喜乐；\BibleQuote{欢乐的}因任何暂时利益而喜乐时，应惧怕永恒的审判；\BibleQuote{购买的}应这样持有财产，不致因过分的爱而紧抓不放；\BibleQuote{享用这世界的}应思虑它正在逝去，并不存留。\BibleQuote{因为这世界的局面正在逝去；但是，}他说，\BibleQuote{我愿意你们无所挂虑}——换句话说：我愿意你们举起心，使它停留在那些不逝去的事物中。然后他接着说：\BibleQuote{没有结婚的，挂虑主的事，想怎样悦乐主；结婚的，挂虑世俗的事，想怎样悦乐妻子。}\newline{}\newline{}\BibleRef{格前 7:32-33} 这样，他在一定程度上解释了他已经说过的话：\BibleQuote{有妻子的要像没有一样。}因为那些有妻子却这样挂虑主的事，想怎样悦乐主，而不挂虑世俗的事以悦乐妻子的人，实际上就像没有妻子一样。当妻子也有这样的心态时，这一点更容易实现，因为她们悦乐丈夫不仅是因为他们富有、地位高、出身高贵、性情温和，而是因为他们是信徒，是虔诚的人，是贞洁的人，是善人。\par

% structure: p0033 source=h2 target=subsection reason=relative-heading-rank; base=h2

\subsection{第十六章 [第十四]——在已婚者身上，某种程度的不节制是应被容忍的；仅为满足情欲而使用婚姻并非无罪，但因婚姻关系，此罪是小罪。}

但在已婚者身上，既然这些事是可欲求、可赞美的，那么另一些事也应被容忍，以免陷入可罚的罪恶中，即奸淫和通奸。为了避免这种恶，即使是那些不以生育为目的、而是为压制性欲服务的夫妻拥抱，也被允许到可以宽恕的程度，尽管不是出于诫命的命令：\newline{}\newline{}\BibleRef{格前 7:6} 而且夫妻被告诫不要互相亏负，免得撒殚因他们的不节制而诱惑他们。\newline{}\newline{}\BibleRef{格前 7:5} 因为经上这样说：\BibleQuote{丈夫对妻子该尽他应尽的义务，妻子对丈夫也是如此。妻子对自己的身体没有主权，而是丈夫有；同样，丈夫对自己的身体也没有主权，而是妻子有。你们不要彼此亏负，除非两厢情愿，暂时分房，为专务祈祷；但事后仍要同房，免得撒殚因你们的不节制而诱惑你们。我说这话，是出于宽容，并不是出于命令。}\newline{}\newline{}\BibleRef{格前 7:3-6} 现在，在必须给予宽容的情况下，绝不能争辩说没有某种程度的罪。然而，既然为了生育子女而结合——这必须被承认为婚姻的真正目的——是无罪的，那么宗徒允许宽容的又是什么呢？无非是：已婚者若没有节制的恩赐，可以彼此要求肉体的债——这不是出于生育的意愿，而是出于情欲的快乐。这种满足并不因婚姻而招致罪责，而是因婚姻而得到宽容。因此，这应被算作婚姻的称赞之一：即婚姻本身使那些本质上不属于它的事成为可宽恕的。因为服务于情欲要求的婚姻拥抱，是这样进行的，并不妨碍生育，即婚姻的目的与目标。\par

% structure: p0035 source=h2 target=subsection reason=relative-heading-rank; base=h2

\subsection{第十七章 [第十五]——在婚姻的使用中，什么是无罪的？什么是附带小罪的？什么是附带大罪的？}

然而，已婚者只为求子而交合，这是另一回事，并非有罪；他们在同居中渴望肉欲快感，但仅限于配偶，则涉及小罪。因为尽管生育后裔并非交合的动机，他们仍不试图以邪恶的欲望或邪恶的器具阻止这种生育。凡采取这些做法的人，虽被称为配偶之名，实则并非配偶；他们不保留真正婚姻的任何痕迹，而是以这光荣的称号作为罪恶行径的掩饰。他们甚至已进展到遗弃自己不愿生下的子女。他们憎恶养育和保留那些他们害怕生下的孩子。这种对勉强生下的后代施加的残忍，揭露了他们在暗中所行的罪恶，并将其清楚地拖入光天化日之下。公开的残忍谴责了隐藏的罪恶。有时，这种好色的残忍，或如果你愿意，残忍的好色，甚至采用如此极端的方法，使用毒药来确保不孕；或者，如果不成功，则通过某种方式在产前毁灭已受孕的种子，宁愿其后代灭亡而不愿获得生命；或者，如果胚胎已在子宫中发育，应在出生前将其杀死。那么，如果双方同样邪恶，他们就不是夫妻；如果从一开始就有这样的品性，他们就不是通过婚姻结合，而是通过淫乱。但如果两人在此罪恶上并不相同，我大胆宣称：那女人可以说是丈夫的娼妇，或者那男人是妻子的奸夫。\par

% structure: p0037 source=h2 target=subsection reason=relative-heading-rank; base=h2

\subsection{第十八章——节制优于婚姻；但婚姻优于奸淫。}

既然婚姻不可能像原始人类若没有罪在先可能有的那样，但它仍可能像古代圣祖们的婚姻那样，即那引起羞愧的\OriginalTerm{肉欲}{carnal concupiscence}（它在堕落前的乐园里不存在，堕落之后也不被允许留在那里），虽然必然成为这死亡之身的一部分，却并不服从它，而只是在被迫时为了帮助生育子女的唯一目的而服从其功能；否则，由于现今时期（正如我们已说过的）是远离婚床的时期，因此并不必然要求行使上述功能，既然现在万民都如此丰盛地贡献于后代——这些后代将获得属灵的生养，那么优良节制的祝福便有更大的空间。\BibleQuote{能领悟的，就领悟吧！}\BibleRef{玛 19:12}他，然而，不能领悟的人，\BibleQuote{即使他结婚，也不是罪；}\BibleRef{格前 7:28}\BibleQuote{如果女人没有节制的恩赐，就让她也结婚吧。}\BibleRef{格前 7:9}\BibleQuote{男人不接触女人，固然是好。}\BibleRef{格前 7:1}\BibleQuote{不是所有人都能领受这话，只有那些得了恩赐的人才能，}\BibleRef{玛 19:9}\BibleQuote{为了避免奸淫，男人当各有自己的妻子，女人当各有自己的丈夫。}\BibleRef{格前 7:2}因此，不节制的软弱便因婚姻的尊贵身份而免于堕入放荡的毁灭。现在，宗徒论到妇女所说的话：\BibleQuote{所以我愿意年轻的女子结婚}\BibleRef{弟前 5:14}也适用于男性：我愿意年轻的男子娶妻，这样便使两性同样\BibleQuote{生育儿女，做母亲、治理家务，不给敌人以诽谤的机会。}\BibleRef{弟前 5:14}\par

% structure: p0039 source=h2 target=subsection reason=relative-heading-rank; base=h2

\subsection{第十九章——婚姻的祝福。}

在婚姻中，让我们爱这些婚姻的恩惠——子女、忠信、\OriginalTerm{圣事性结合}{sacramental bond}。子女，不仅是为了出生，更是为了重生；因为若不为生命而重生，生来就是为了惩罚。忠信，并非像不信者之间彼此保持的那种，出于对肉体的炽爱。因为无论多么不虔敬的丈夫，谁会喜欢奸淫的妻子？或无论多么不虔敬的妻子，谁会喜欢奸淫的丈夫？这确实是婚姻中自然的善，虽是属肉的善。但基督的肢体不应为了自己，而应为了他的配偶害怕奸淫；并应希望从基督那里得到他对配偶所表现之忠信的赏报。再者，圣事性结合，既不因离婚也不因奸淫而丧失，应由夫妻以和谐与贞洁守护。因为唯独它是不孕的婚姻也藉着虔敬的法律所保守的，即使为了婚姻的一切果实希望都已丧失。愿称赞婚姻者赞扬婚姻中这些婚姻的恩惠。然而，\OriginalTerm{肉欲}{carnal concupiscence}不可归咎于婚姻：它只是在婚姻中被容忍。它不是来自婚姻本质的善，而是原罪偶然带来的恶。\par

% structure: p0041 source=h2 target=subsection reason=relative-heading-rank; base=h2

\subsection{第二十章——为何从神圣婚姻生出的是义怒之子。}

这就是为什么即使天主子女的合法婚姻所生的孩子，也不是天主的孩子，而是世界的孩子；因为那些生育者，如果他们已经\OriginalTerm{重生}{regeneration}，所生的孩子不是作为天主的孩子，而是仍然作为世界的孩子。主说：\BibleQuote{今世之子，}\BibleQuote{生养并受生。}因此，由于我们仍然是今世之子，我们的\OriginalTerm{外在的人}{outer man}处于败坏状态；因此，我们的后代出生为现世的孩子；除非他们重生，否则不会成为天主的儿子。然而，既然我们是天主的子女，我们的\OriginalTerm{内在的人}{inner man}得以日日更新。\BibleRef{格后 4:16} 而且，即便是我们的外在的人，也通过重生之洗得以圣化，并获得了将来不朽的望德，因此它被恰当地称为\BibleQuote{天主的殿}。宗徒说：\BibleQuote{你们的身体是住在你们内之圣神的宫殿，这圣神是你们从天主所受的；你们不是属于自己，因为你们是用重价买来的；所以要在你们身上光荣并承载天主。}这一切陈述都是关于我们目前的圣化，但特别是基于他对另一处所讲的盼望：\BibleQuote{我们这些有圣神初果的人，我们自己也在心中叹息，等待着义子名分，就是我们身体得救赎。}\BibleRef{罗 8:23} 如果身体得救赎是所盼望的，如宗徒所说，那么既然是盼望，就仍是望德之事，而非实际拥有。因此宗徒接着说：\BibleQuote{我们得救在于望德；可见的望德不是望德：因为人看见的，为什么还盼望呢？但我们若盼望那看不见的，就要耐心等待。}\BibleRef{罗 8:24} 所以，不是由于我们所等待的，而是由于我们现在所忍受的，我们肉身的孩子才出生。一个有信德的人，当他听到宗徒吩咐人\BibleQuote{爱自己的妻子}\BibleRef{哥 3:19}时，万不可爱他妻子身上那肉身的\OriginalTerm{私欲偏情}{concupiscence}，就是他自己身上也不该爱的；他若听从另一位宗徒的话，就可以知道：\BibleQuote{不要爱世界，也不可爱世界上的事。谁若爱世界，父的爱就不在他内。因为凡世界上的事，肉身的欲望、眼目的欲望、人生的骄傲，都不是出于父，而是出于世界。世界和它的欲望都要逝去；但承行天主旨意的人，却永远长存，如同天主永远长存一样。}\par

% structure: p0043 source=h2 target=subsection reason=relative-heading-rank; base=h2

\subsection{第二十一章〔第十九〕—— 因此，罪人由义人父母所生，正如野橄榄从橄榄树生出一样。}

因此，那由肉身贪欲所生的，实在是生于世界，而非生于天主；但当它由水和圣神重生时，便生于天主。这私欲偏情的罪债，只有重生才能赦免，正如自然的生育背负它一样。那么，那被生的必须重生，以便同样地，既然不能不然，所背负的得以赦免。无疑，这是非常奇妙的：在父母身上被赦免的，却在子女身上被背负；但事实确是如此。为使这奥秘的真理——不信者既看不见也不相信——获得一些有形的证据来支持它，天主以祂的眷顾在某些树中提供了例子。因为我们为何不认为，正是为此目的，野橄榄从橄榄树生出？难道不的确可信吗？在为人所用而受造的事物中，造物主提供并指定了什么来提供一个适用于人类的教育性例子？那么，这是奇妙的事：那些本身已因恩宠从罪的束缚中获得释放的人，怎能仍然生出那些被同一锁链捆绑、并需要同样解绑过程的人？是的；我们承认这奇妙的事实。但是，野橄榄树的胚胎潜在存在于真橄榄的种子里，若不是通过实验和观察证实，谁会相信呢？因此，正如野橄榄从野橄榄的种子长出，从真橄榄的种子长出的也只是野橄榄，尽管野橄榄与橄榄之间存在着极大的差异；同样，在肉身中所生的，无论是于罪人或于义人，两者都是罪人，尽管罪人与义人之间存在着巨大的区别。那被生的，在行为上还不是罪人，从出生来说还是新的；但在罪债上他却是老的。作为来自造物主的人，他是毁灭者的俘虏，需要一位救赎者。然而，困难在于，当父母本身已先前从被掳中获得救赎时，被掳的状态怎能临到后代？解开这个棘手的点，或在一篇系统的论述中解释它，并非易事；因此不信者拒绝接受它为真；正如在关于野橄榄和橄榄的那个比喻中，我们给出的解释，任何理由都能轻易找到，或者清楚地解释为什么同样的嫩芽能从如此不同的树干长出。然而，任何愿意做实验的人都能发现这其中的真理。那就让它成为一个好例子，暗示相信那无法眼见之事。\par

% structure: p0045 source=h2 target=subsection reason=relative-heading-rank; base=h2

\subsection{第二十二章〔第二十〕—— 即使婴儿，未受洗时，也在魔鬼权下；针对婴儿的驱魔，及弃绝魔鬼。}

基督信仰坚定地宣告，我们这些新异端者开始否认的，即那些藉重生之洗而得洁净的人已从魔鬼的权下被赎回，以及那些尚未藉此重生被赎回的人仍在魔鬼的权下被掳，即使他们是已被赎回者的婴孩儿女，除非他们藉同一基督的恩宠被赎回。因为无异，天主的祝福适用于人类生命的每一阶段，宗徒论及祂时说：\Quote{祂拯救了我们脱离黑暗的权势，把我们迁到祂爱子的国里。} \BibleRef{哥 1:13} 因此，婴孩受洗时，从黑暗的权势（魔鬼是这权势的君王），换言之，从魔鬼及其天使的权势下被救出；凡否认此点的，被教会圣事的真理所定罪，这圣事在基督的教会中，任何异端新说都不得毁坏或改变，只要神圣的元首支配并帮助祂所拥有的整个身体——无论小的还是大的。因此，真实无伪，婴孩被驱魔，他们藉着带他们来受洗者的心和口（自己不能这样做）弃绝魔鬼，为使他们从黑暗的权势下被救出，并被迁到他们主的国里。那么，在他们里面，是什么使他们留在魔鬼的权下，直到藉基督的圣洗圣事被救出？我问，除了罪还有什么？确实，魔鬼没有找到别的，使他能将那善的造物主所造善的人性置于自己控制之下。但婴孩自出生以来没有犯过自己的罪。因此只有原罪留下，使他们被掳于魔鬼的权下，直到藉重生之洗和基督的血被赎回，进入他们救赎主的国——那掳掠者的权势被挫败，并且他们得权能成为\Quote{天主的子女}，而非今世之子。 \BibleRef{若 1:12}\par

% structure: p0047 source=h2 target=subsection reason=relative-heading-rank; base=h2

\subsection{第二十三章 — 罪并非从婚姻的善性中产生；婚姻的圣事在基督与教会的情形中是伟大的，在一男一女的情形中则很小。}

如果现在我们（可以这么说）询问我们常提到的那些婚姻的善，并询问罪如何可能从它们传播给婴孩，我们将从第一个善——繁衍后代的工作——得到这样的回答：\Quote{我的幸福在乐园中本会更大，如果罪没有被犯下。因为全能天主的祝福属于我：\InnerQuote{你们要生育繁殖。} \BibleRef{创 1:29} 为完成这一善工，创造了适合各性别的不同肢体；这些肢体当然在罪之前就已存在，但它们并非羞耻的对象。} 第二个善——贞洁的信实——将这样回答：\Quote{如果罪没有被犯下，在乐园中有什么比我自己更安全呢？那时没有我自己的私欲来刺激我，也没有他人的私欲来诱惑我。} 然后第三个善——婚姻的圣事性结合——将这样回答：\Quote{在罪进入之前，在乐园中说了那句话：\InnerQuote{人应离开自己的父母，依附自己的妻子，两人成为一体。} \BibleRef{创 2:24} } 宗徒将此应用于基督与教会的情形，并称之为\Quote{伟大的圣事}。那么，在基督与教会中伟大的，在每一对已婚夫妻的情形中则很小，但即便如此，它仍是不可分离结合的圣事。在这三种婚姻的祝福中，有什么能使罪的束缚传给后代呢？绝对没有。可以肯定，婚姻的善完全包含在这些祝福中；甚至如今，好的婚姻也由这些同样的祝福构成。\par

% structure: p0049 source=h2 target=subsection reason=relative-heading-rank; base=h2

\subsection{第二十四章 — 私欲与羞耻源自罪；罪的法律；犬儒学派的无耻。}

但如果同样询问\OriginalTerm{肉身的私欲}{concupiscentia carnis}，如今那些一度没有羞耻的行为是如何带来羞耻的，她的回答难道不是，她只在罪之后才开始存在于人的肢体中？[XXII.] 因此，宗徒将她的影响称为\OriginalTerm{罪的法律}{lex peccati} \BibleRef{罗 7:23}，因为她使人屈服于自己，当人不愿服从自己的天主时；正是她使第一对已婚夫妇在那一刻感到羞耻，当他们遮住自己的腰部；正如所有人仍然感到羞耻，寻找隐秘处所同居，甚至不敢让那些他们自己如此生下的孩子看到自己所做的事。正是针对这种自然羞耻的谦逊，犬儒派哲学家，在令人惊愕的无耻之误中，如此激烈地挣扎：他们认为夫妻的交合既然是合法且可敬的，就应当在公开场合进行。如此厚颜无耻的猥亵活该被称为狗；因此他们被称为\Quote{\Emph{\OriginalTerm{犬儒}{Cynics}}}。\par

% structure: p0051 source=h2 target=subsection reason=relative-heading-rank; base=h2

\subsection{第二十五章 — 在重生者中未经同意的私欲并非罪；私欲在何种意义上被称为罪。}

现在，这私欲偏情，这居于我们肢体中的罪律——正义的律法禁止我们效忠于它，藉宗徒的话说：\Quote{所以，不要让罪在你们可朽的身体内为王，使你们顺从它的情欲；也不要把你们的肢体交给罪，作不义的工具：} \BibleRef{罗 6:12}——这私欲偏情，我说，它唯独藉重生的圣事才能被洁净，无疑地，它藉自然的生育，将罪的锁链传给人的后裔，除非他们自己藉重生获得解脱。然而，在重生者身上，只要他们不为私欲偏情所动而行非法之事，且不让主宰的理性运用肢体去犯这样的行为，这私欲偏情就不再是罪了。因此，若在一处所吩咐的\Quote{不可贪恋}没有遵守，至少另一处所命令的\Quote{不可随从你的私欲偏情}得以遵守。\BibleRef{德 18:30} 然而，由于它凭某种说法被称为罪——因为它从罪而来，且当它占上风时就产生罪——它的罪责在属血气的人身上得势；但这罪责，因基督的恩宠藉诸罪的赦免，在重生者身上不被容许得势，只要他不顺从它每当它催促他作恶时的鼓动。我说，它既然从罪而来，就称为罪，虽然在重生者身上实际上不是罪；它被赋予这个名称，就像舌头所产生的言语本身被称为\Quote{\Emph{舌头}}；又如\Quote{\Emph{手}}这个词被用来指手所产生的书写。同样，私欲偏情被称为罪，因为它当战胜意志时就产生罪：就如寒冷和霜冻被冠以\Quote{迟钝}的称号；并非因迟钝而来，而是产生迟钝；实际上使我们麻木。\par

% structure: p0053 source=h2 target=subsection reason=relative-heading-rank; base=h2

\subsection{第二十六章 凡因私欲偏情而生者，无不因罪而理所应当地受制于魔鬼；魔鬼应受比人更重的惩罚。}

魔鬼加诸人类身上的这个创伤，迫使一切因此创伤而生的事物都处于魔鬼的权势之下，如同他正当地从自己的树上摘下果实一样。并非因为人的本性——它唯独来自天主——出自魔鬼，而是因为唯独罪——它不是来自天主——出自魔鬼。因为人的本性并非因其自身而受定罪，它是天主的化工，因此是可赞美的；而是因其可定罪的败坏，这败坏了它。正是因这定罪，它受制于魔鬼，而魔鬼也处于同样可罚的状态。因为魔鬼本身是一个不洁之神：作为神体，他原是善的；但因不洁而成为恶的；因为他本性上是神体，因败坏而成为不洁的。在这两者中，一者来自天主，另一者来自他自己。因此，他控制人——无论老年人还是婴儿——不是因为他们是人，而是因为他们被玷污。所以，谁若惊讶于天主的受造物受制于魔鬼，就当停止这种惊讶。因为天主的受造物受制于另一个天主的受造物，较小的受制于较大的，人受制于天使；而这并非由于本性，而是由于本性的败坏：统治者被玷污，被统治者也被玷污。这一切都是那古老败坏之种的果实，这败坏他已种植在人里面；他自己注定要在最后审判中受更重的惩罚，因为他是更被玷污的；但同时，那些在定罪中将承受较轻负担的人，也作为罪之君王与创始者的他之下，因为除罪以外没有别的定罪原因。\par

% structure: p0055 source=h2 target=subsection reason=relative-heading-rank; base=h2

\subsection{第二十七章 ［XXIV］——原罪藉情欲传递；夫妻间的小罪；肉身的私欲偏情，罪之女与罪之母。}

因此，魔鬼使婴儿成为有罪的——他们不是生自那使婚姻为善的善，而是生自私欲偏情的恶，这恶婚姻虽正确使用，但连婚姻也有理由对此感到羞耻。婚姻本身在一切正当属于它的善中\Quote{在众人中都是可尊敬的} \BibleRef{希 13:4}；但即使它有\Quote{没有玷污的床第}（不仅远离邪淫和奸淫——这些是可定罪的耻辱——也远离任何过度的同居行为，这些不来自对子女的强烈愿望，而来自难以驾驭的享乐情欲，是夫妻间的小罪），然而，只要进入实际生育过程，那合法且可敬的拥抱本身，若没有情欲的热火，就不能实现，从而完成那属于理性使用而非情欲的事。现在，这热火，无论跟随还是先于意志，都以一种自身的力量激动那些不能单纯由意志运动的肢体，借此表明它不是那命令它的意志的仆人，而是那违背它的意志的惩罚。它还表明，它必须不是由自由选择来激发，而是由某种诱惑性的刺激，并且正是因此它产生羞耻。这就是肉身的私欲偏情，它在重生者身上虽不再算为罪，但除了从罪而来外，绝不出现在本性中。它可以说是罪之女；而每当它同意去行可耻之事时，它就又成为许多罪之母。从此私欲偏情而来的一切，由自然生育进入存在，都被原罪所束缚，除非它确实在祂内重生——祂是童贞玛利亚不藉此私欲偏情而怀孕的。因此，当祂屈尊以肉身诞生时，唯独祂是无罪而生的。\par

% structure: p0057 source=h2 target=subsection reason=relative-heading-rank; base=h2

\subsection{第二十八章 ［XXV］——私欲偏情在圣洗后仍存留，正如虚弱在疾病痊愈后仍存留；私欲偏情在年岁增长者中减少，在不能自制者中增加。}

如果发生这样的问题：\OriginalTerm{肉性私欲}{concupiscentia carnis}如何在重生者身上存留——在他们身上已经获得了一切罪的赦免——因为人类的生育是通过它进行的，即使是一个已受洗的父母的肉身后代出生；或者无论如何，如果在一个已受洗的父母身上，这私欲可能存在而不成为罪，那么为什么在儿女身上同样的私欲就成为罪呢？——答案是这样的：肉性私欲确实在圣洗中得到了赦免，但不是使之消失，而是使之不归算为罪。虽然它的罪责现在被除去，但它仍然存留，直到我们整个的软弱借着内在的人日益更新的进展而被医治，那时我们外在的人最终将穿上不朽。\BibleRef{格前 15:53}然而，它并非作为实体，如身体或灵魂那样存留，它不过是一种邪恶品质的倾向，例如衰弱。当然，没有任何东西遗留需要被赦免，正如圣经所说：\BibleQuote{主赦免我们的一切罪过。}但直到同一段经文紧接着发生的事：\BibleQuote{祂医治你的一切病苦，救赎你的性命脱离死亡，}这肉性私欲仍存留在这死亡的身体中。现在，我们被劝诫不要顺从它邪恶的欲望去作恶：\Quote{不要让罪在你们必死的身体上为王。}\BibleRef{罗 6:12}然而，这私欲在节制的人和年岁增长的人中日益减少，尤其在老年临近时。但是，那向它屈从邪恶服务的人，会获得如此强大的力量，以至于即使他所有的肢体因年龄而衰败，他身体那些特殊的部分无法履行其正常功能，他也绝不停止在日益增长的卑鄙无耻欲望的激情中放纵。\par

% structure: p0059 source=h2 target=subsection reason=relative-heading-rank; base=h2

\subsection{第二十九章〔XXVI〕——论私欲偏情在已受洗者身上如何行动上残留，但罪责已消除。}

那么，对于那些在基督内重生的人，当他们获得一切罪的全部赦免时，当然也必须赦免仍然内住的私欲偏情的罪责，以便（如我所说）它不被归算为他们的罪。因为即使在那些不能持久、一犯即过的罪的情况下，罪责却是持久的，并且（若得不到赦免）将永远存留；所以当私欲偏情被赦免时，它的罪责也被除去。因为\Quote{没有罪}的意思就是不被视为有罪。如果一个人（例如）犯了奸淫，即使他不再重复这罪，他仍被视为有罪，直到罪过本身得到赦免。因此，他的罪仍然存留，尽管他罪行的具体行为已不复存在，因为它已随犯下的时间而消逝。因为如果停止犯罪等同于没有罪，那么圣经只需简单地警告我们：\BibleQuote{我儿，你犯了罪吗？不要再犯。}\BibleRef{德 21:1}这就足够了。然而，这还不够，因为它接着说：\BibleQuote{求你赦免你从前的罪过。}\BibleRef{德 21:1}因此，罪若不得到赦免，就仍然存留。但是，如果罪已消逝，它们如何存留呢？只能这样：它们在\Emph{行为}上已经消逝，但在\Emph{罪责}上却持久。相反，则可能发生这样的情况：一件事可能在行为上存留，但在罪责上消逝。\par

% structure: p0061 source=h2 target=subsection reason=relative-heading-rank; base=h2

\subsection{第三十章〔XXVII〕——论私欲偏情的邪恶欲望；我们应当期望它们不存在。}

肉性私欲在某种程度上是活跃的，即使它既不表现出内心的同意（那是它的统治宝座），也不表现出那些用作武器来完成其意图的肢体。但在此行动中，它产生什么，除了它邪恶可耻的欲望？因为如果这些欲望是善而合法的，宗徒就不会禁止顺从它们，说：\BibleQuote{所以不要让罪在你们必死的身体上为王，以使你们顺从它的情欲。}\BibleRef{罗 6:12}他没有说，你们有它的情欲，而是说\BibleQuote{你们顺从它的情欲}；为使我们（由于欲望在不同人中大小不同，视各人在内在之人的更新上进展如何）可以保持圣洁和贞洁的战斗，目的是拒绝顺从这些情欲。然而，我们的愿望应该是这些欲望根本不存在，即使这样的愿望在这死亡的身体中不可能实现。这就是为什么同一宗徒在另一段落中，仿佛以自己个人的身份对我们说，给我们这个教训：\BibleQuote{我所愿意的，}他说，\BibleQuote{我不做；我所憎恨的，我反而去做。}\BibleRef{罗 7:15}简言之，\Quote{我贪恋。}因为他不想做这事，以便在各方面都完美。\BibleQuote{如果，我做了我所不愿意的，}他继续说，\BibleQuote{我就同意法律是善的。}\BibleRef{罗 7:16}因为法律也不愿意我所不愿意的。因为法律不愿意我有私欲偏情，它说：\Quote{不可贪恋；}而我也不愿意怀有如此邪恶的欲望。因此，在这里，法律的意志和我自己的意志完全一致。但是因为他不想贪恋，却仍然贪恋，而且尽管如此，他绝不服从这私欲偏情以致同意它，他立即加上这些话：\BibleQuote{那么，现在不是我做的，而是住在我内的罪做的。}\BibleRef{罗 7:17}\par

% structure: p0063 source=h2 target=subsection reason=relative-heading-rank; base=h2

\subsection{第三十一章〔XXVIII〕——谁能说\Quote{不是我做}？}

一个人，若同意自己肉体的贪欲，并在心中决定行其所欲，且着手去做，却自以为仍有权利说：\Quote{那不是我在做}，即使他因同意邪恶欲望而恨恶并厌恶自己，那就大错特错了。这在他身上是同时发生的：他本人恨恶那事，因为他知道那是邪恶的；但他仍去做，因为他决意去做。此外，如果他还进一步去做圣经所禁止的事，当圣经说：\BibleQuote{不要将你们的肢体献给罪恶，作不义的武器}\BibleRef{罗 6:13}，并且以身体的行为完成他在心中决意要做的事，然后说：\BibleQuote{不是我做的，而是住在我里面的罪做的}\BibleRef{罗 7:17}，因为他对自己决心并完成那事感到不满——他就如此大错，以致不认识自己。因为，当他整个自我都是他自己——他的理智决定，他的身体执行他自己的意图——他却以为自己不再是原来的他了！[XXIX.]因此，只有那人才说真话，他说：\BibleQuote{不再是我做的，而是住在我里面的罪做的}，即那人只感受到贪欲，却没有用心中的同意去决定做它，也没有用身体的服从来完成它。\par

% structure: p0065 source=h2 target=subsection reason=relative-heading-rank; base=h2

\subsection{第32章——何时完美行善}

宗徒随后加上这些话：\BibleQuote{因为我知道，在我内，即在我肉体内，没有善；因为愿意善是与我同在，但如何完成善，我却找不到。}\BibleRef{罗 7:18} 这是说，善事尚未完美，当没有邪恶欲望时，如同邪恶在邪恶欲望被服从时被完美成全一样。但当它们存在却未服从时，既没有行恶，因为没有顺服它们；也没有行善，因为它们无效的存在。那里是一个中间状态：善在一定程度上被实现，因为邪恶贪欲没有获得同意；恶在一定程度上存留，因为贪欲存在。这解释了宗徒精确的措辞。他没有说\Emph{做}善不与他同在，而是说\Quote{如何完成它}。事实上，一个人做了大量的善，当他遵行圣经所吩咐的：\BibleQuote{不要随从你的私欲}\BibleRef{德 18:30}；然而他仍达不到完美，因为他未能遵守那大诫命：\BibleQuote{不可贪恋}\BibleRef{出 20:7}。法律说：\BibleQuote{不可贪恋}，为使我们发现自己处于此病态时，可寻求恩宠的良药，并借着那诫命得知在我们现世可朽坏的状态中应朝哪个方向努力，以及在未来的不死中能达到何等高度。因为除非完美能在某处达到，这条诫命永远不会赐给我们。\par

% structure: p0067 source=h2 target=subsection reason=relative-heading-rank; base=h2

\subsection{第33章[XXX.]——真正的自由来自甘愿喜乐意愿天主的法律}

宗徒接着重复他先前的陈述，以更充分地推荐其要旨：\BibleQuote{因为我所愿意的善，我不做；我所不愿意的恶，我倒去做。若我所不愿意的，我倒去做，就不再是我做的，而是住在我里面的罪做的。}接着是：\BibleQuote{因此我发现，当我愿意行善时，法律是善的；因为恶与我同在。}换句话说，我发现，当我愿意做法律所要求的事时，法律对我有益；因为并不是法律本身（它说：\Quote{不可贪恋})中有什么恶；不，是恶与我同在，即使我愿意，我仍有贪欲。\Quote{因为，}他补充说，\BibleQuote{我按内在的人欣悦天主的法律；但我看见在我肢体中另有一条法律，与我的理智法律交战，并把我俘掳，隶属于那在我肢体中的罪的法律。}\BibleRef{罗 7:22} 这种按内在的人对天主法律的欣悦，来自天主的强大恩宠；借此我们的内在之人日新又新\BibleRef{格后 4:16}，因为我们借此坚持进步。在其中没有畏惧所带来的痛苦，而是喜爱所带来的愉悦和满足。我们在那里是真正自由的，在那里我们没有不情愿的喜乐。\par

% structure: p0069 source=h2 target=subsection reason=relative-heading-rank; base=h2

\subsection{第34章——贪欲如何俘掳了宗徒；罪的法律对宗徒是什么}

于是，这句话：\BibleQuote{我看见在我肢体中另有一条法律，与我的理智法律交战}指的是我们现在正谈论的贪欲本身——我们肉体中罪的法律。但当他说：\BibleQuote{并把我俘掳，隶属于罪的法律}，即隶属于它自身，\Quote{那在我肢体中的}，他要么意指\Quote{把我俘掳} 在\Emph{努力使我}被俘掳的意义上，即催促我赞同并实现邪恶欲望；要么（这并不引起争论）指按照肉体俘虏我，如果肉体没有被肉性的贪欲（他称之为罪的法律）所占据，当然就没有不正当的欲望——我们的理智不应服从的欲望——来激起并扰乱。然而，宗徒没有说：把\Emph{我的肉体}俘掳，而是说：\Quote{把我俘掳}，这引导我们寻找另一种含义，并将\Quote{把我俘掳}理解为好像他说：努力使我被俘掳。但为什么他不可以说：\Quote{把我俘掳}，同时意指我们的肉体？难道关于耶稣，当他的肉体在坟墓中找不到时，不是有人这样说：\BibleQuote{他们把主拿走了，我不知道他们把祂放在哪里}\BibleRef{若 20:2}？那么玛利亚的问题就不恰当吗，因为她说了：\BibleQuote{我的主}，而不是\Quote{我主的身体}或\Quote{肉体}？\par

% structure: p0071 source=h2 target=subsection reason=relative-heading-rank; base=h2

\subsection{第35章[XXXI.]——肉体，肉性的爱欲}

但我们在宗徒自己的话里，稍前一点，有一个足够清楚的证明，证明当他所说的\Quote{把\Emph{我}掳去}时，他可能是指他的肉体。因为宣告了\BibleQuote{我知道在\Emph{我}之内没有善存在}之后，他立刻加上一句解释的话说：\BibleQuote{就是说，在我的肉体之内。} \BibleRef{罗 7:18} 那么，就是那肉体——里面没有善居住——被掳去服从罪恶的法律。现在他称那肉体为某种病态的肉身贪恋，而非我们身体构造的单纯形状，其肢体不应被用作罪恶的武器——即不应被用作那俘获我们肉体的私欲偏情本身。就我们这实际的身体实体和本性而言，在所有忠信的人中，无论他们生活在婚姻中还是节欲中，它已经是天主的宫殿。然而，如果我们的肉体绝对没有任何部分被俘虏，甚至被魔鬼俘虏（虽然因着罪的赦免，罪恶已不被归咎，这正被恰当地称为罪恶的法律），那么，如果在罪恶的法律之下，即在其自身的私欲偏情之下，我们的肉体没有在一定程度上被俘虏，宗徒所说的我们\BibleQuote{等待义子的身分，就是我们身体的救赎}又怎能是真的呢？ \BibleRef{罗 8:23} 因此，既然如今我们仍在等待身体的救赎，那么在我们之内也仍有一些东西被俘虏于罪恶的法律之下。因此他喊道：\BibleQuote{我真不幸！谁能救我脱离这死亡的身体？天主的恩宠，借着我们的主耶稣基督。} \BibleRef{罗 7:24} 我们如何理解这样的话，除了说我们这正在腐朽的身体重压着我们的灵魂？因此，当我们这身体本身在未来以不朽的状态恢复给我们时，我们将彻底脱离这死亡的身体；但对那些将要复活而受审判的人，却没有这样的解救。那么，这死亡的身体被理解为导致我们肢体中有另一条法律与心智的法律交战，只要肉体与心神交战——但并没有征服心神，因为心神方面也有相反肉体的贪欲。 \BibleRef{迦 5:17} 这样，尽管罪恶的法律部分地俘虏了肉体（由此肉体抵抗心智的法律），但它在我们的身体中并没有绝对的统治权，尽管身体处于必死的状态，因为它拒绝顺从它的欲望。 \BibleRef{罗 6:12} 因为在敌对军队之间，当有激烈的冲突时，甚至战斗中较弱的一方通常也会持有所俘获之物；虽然以某种方式，我们的肉体中有一些东西被保持在罪恶的法律之下，但它仍有望获得救赎：那时将不再存留一丝这种腐朽的私欲偏情；但我们的肉体，治好了那疾病的瘟疫，完全披上不朽，将在永恒的幸福中永远生活。\par

% structure: p0073 source=h2 target=subsection reason=relative-heading-rank; base=h2

\subsection{第三十六章——即使我们现在仍有私欲偏情，我们在基督内仍可得安全。}

但宗徒继续论述说：\BibleQuote{这样看来，我本人以心思服事天主的法律，而以肉体服事罪恶的法律。} \BibleRef{罗 7:25} 这必须如此理解：\BibleQuote{我以心思服事天主的法律，}是因为我拒绝同意罪恶的法律；\BibleQuote{我以肉体}服事\BibleQuote{罪恶的法律，}是因为我有罪的欲望，我尚未完全从这些欲望中解放，虽然我并不表示同意。然后让我们仔细观察他在上述一切之后所说的话：\BibleQuote{因此，现今在基督耶稣内的，已没有定罪。} \BibleRef{罗 8:1} 他说，甚至是\Emph{现今}，当在我肢体中的法律与我的心智的法律争战，并在死亡的身体中保留一些俘虏时，在基督耶稣内的已没有定罪。并且听原因：\BibleQuote{因为生命之神的法律在基督耶稣内}，他说，\BibleQuote{使我脱离了罪恶与死亡的法律。} \BibleRef{罗 8:2} 它如何使我自由？除非是因着赦免我所有的罪，它取消了罪的判决；这样，虽然它仍然存在，只是每日逐渐减少，但它不再被归咎为罪？\par

% structure: p0075 source=h2 target=subsection reason=relative-heading-rank; base=h2

\subsection{第三十七章[三十二]——未受洗婴儿中罪恶的法律及其罪责。因亚当的罪，人类变成了一棵\Quote{野橄榄树}。}

直到这罪之赦在后裔身上发生为止，他们内在的罪律确实被归为罪；换言之，该律法附带罪的判决，使他们成为永罚的债务人。因为父母传给血肉子女的是他自己血肉出生的状况，而非他属神重生的状况。他虽在血肉中出生，但在他的罪过得赦后，这并不妨碍他结出果实，然而这仍像隐藏于橄榄种子中一样；甚至因他的罪得赦，这丝毫未损及油——即通俗而言，他在基督之后所活的性命，是\Quote{因信德而成义的人，} \BibleRef{罗 1:17} 在基督之后，基督的名字本身即源于油，即源于傅油。然而，在重生的父母身上，如同在纯橄榄的种子中，那已获赦免而无罪责的部分，仍然毫无疑问地保留在他的未重生的后裔身上，如同在野橄榄中，带有全部罪责，直到这罪责也由同一恩宠获得赦免。当亚当犯罪时，他从那纯橄榄（它没有能生出野橄榄苦果的腐败种子）变成了野橄榄树；而且，因他的罪如此之大，以致他的本性随之变坏，他把整个人类变成了野橄榄的根株。这种改变的效果，正如上文所说，可以在这些树本身看到。每当天主的恩宠将一棵树苗转变为好橄榄时，使得初生的缺陷（即从肉身的贪欲所衍生和沾染的原罪）得以赦免、遮掩、不被归罪，但其中仍然固有那种能生出野橄榄的本性，除非这树苗也藉同一恩宠、由第二次诞生变为好橄榄。\par

% structure: p0077 source=h2 target=subsection reason=relative-heading-rank; base=h2

\subsection{第三十八章 [第三十三]——一切罪的赦免和复活的全然痊愈都必须归于圣洗。每日的洁净。}

因此，这橄榄树是有福的，\Quote{它的过犯得赦免，它的罪恶被遮盖}；它是有福的，\Quote{主不将罪归于它}。但这已获得赦免、遮掩和宣告无罪（甚至直到全然改变为永恒不朽）的，仍然保留一种隐秘的力量，为野而苦的橄榄树提供种子，除非同样的天主耕种也藉赦免、遮掩和宣告无罪来修剪它。然而，当同样的重生（现藉圣洗池生效）清除并治愈人的一切恶直到尽头时，在甚至血肉的种子中也将不再留下任何腐败。藉此，曾经形成血肉心智的同一血肉将成为属神的——不再有那抵抗心神的律法的肉欲，不再发出血肉的种子。因为我们必须如此理解那位我们多次引用的宗徒在别处所说的话：\Quote{基督爱了教会，并为她舍了自己，以便用水藉着言语来圣化并洁净她，好使她能献给自己一个荣耀的教会，没有瑕疵、皱纹或任何这类东西。} \BibleRef{弗 5:25} 我说，必须理解为：藉此重生的洗涤和圣化之言，重生之人的一切恶（无论何种）都被洁净和治愈——不仅是在圣洗中全部得赦免的罪，也包括圣洗后因人的无知和软弱所犯的罪；当然，并非每次犯罪都要重复圣洗，而是藉其唯一一次的施行，使信友在重生前后的一切罪都获得赦免。因为无论是在圣洗之前的忏悔（如果圣洗不随后发生）还是在圣洗之后的忏悔（如果圣洗不先发生）有什么用处呢？不，在主祷文本身（我们每日的洁净）中，念出\Quote{宽免我们的罪债，} \BibleRef{玛 6:12} 的祈求，若非由已受洗的人发出，又有何益处呢？同样，一个人施舍的慷慨与仁慈，无论多大，我问，如果他没有受洗，这些对他罪的赦免有何益处呢？简言之，除了受洗的人，还有谁能得到天国本身的福乐？在那里，教会将没有瑕疵、皱纹或任何这类东西；在那里，将没有可指责的，没有虚幻的；在那里，不仅没有罪的罪责，也没有激起罪的贪欲。\par

% structure: p0079 source=h2 target=subsection reason=relative-heading-rank; base=h2

\subsection{第三十九章 [第三十四]——藉圣洗的圣洁，不仅罪，而且所有种种邪恶都必须被除去。教会尚未从所有瑕疵中得自由。}

因此，不仅所有的罪，而且人的一切种类的疾病，都在藉着那基督徒洗涤的圣洁而被除去，基督以此洁净祂的教会，使她可以献给自己，不是在此世，而是在来世，没有瑕疵、皱纹或任何这类东西。现在有些人主张教会如今就是这样，而他们却在其中。那么，既然他们承认自己有一些罪，如果他们在这一点上说真话（当然他们确实如此，因为他们没有脱离罪），那么教会就在他们身上有\Quote{瑕疵}；而如果他们在宣认中说谎（如同心怀二意地说话），那么教会就在他们身上有\Quote{皱纹}。然而，如果他们断言是他們自己，而非教会，有这一切，那么他们就等于承认自己不是她的成员，不属于她的身体，以致他们甚至被自己的宣认所定断。\par

% structure: p0081 source=h2 target=subsection reason=relative-heading-rank; base=h2

\subsection{第四十章 [第三十五]——藉圣盎博罗削的权威驳斥白拉奇主义者，他们引用他以证明肉身的欲望是自然的美善。}

然而，对于这肉身的私欲偏情，我们在这冗长的讨论中力求将其与婚姻的善准确区分。我们这样做是因为那些现代异端者，每当私欲偏情受到谴责时，他们就吹毛求疵，仿佛这是对婚姻的谴责。他们的目的是赞美私欲偏情为一种自然的善，以便为他们自己有害的教义辩护，即声称通过它而出生的那些人并不染原罪。如今，圣\Person{盎博罗削}，\Place{米兰}主教，我曾藉他的司铎职务领受了重生的圣洗圣事，就此问题简述如下：当他诠释\Person{依撒意亚}先知时，他从先知那里得出基督降生成人的道理：圣\Person{盎博罗削}主教说：\Quote{因此，祂在各方面都像人一样受过试探\BibleRef{希 4:15}，又带着人的形状承受一切；但祂既由圣神而生，就使自己远离罪恶。因为人人都是说谎者，除天主独一外，没有无罪者。所以，历来坚定地认为，显然没有一个人是从丈夫和妻子，即通过他们个人的结合，而无罪的。那无罪者也免于这类受孕。}那么，圣\Person{盎博罗削}在此正确阐明的教义中谴责了什么？——是婚姻的善，还是这些异端者（尽管他们当时尚未登场）的卑劣见解？我以为值得引用这一见证，因为\Person{白拉奇}这样称赞\Person{盎博罗削}说：\Quote{圣\Person{盎博罗削}主教，在其著作中罗马信仰比其他任何地方都更清楚地被陈述，他在拉丁作家中如同一朵美丽的花绽放。他的忠诚和对圣经含义极其纯洁的理解，没有任何对手敢于质疑。}我希望他后悔持有与\Person{盎博罗削}相反的意见，而不是他给予那位圣人的赞美。\par

那么，这就是我的书，由于它冗长的篇幅和艰深的主题，无论对我创作还是对你阅读，都是一件麻烦事——你只能在那些你能够（或至少，我认为你能够）找到闲暇的零碎时间里阅读。尽管我在教会职务中，靠着天主的助佑，着实费了不少心力才写成此书，但我本不应在您公务繁忙之际打扰您，若非一位与您深交的虔诚人士告诉我说，您如此喜爱阅读，以致常常整夜不眠，花费宝贵时光沉浸于您所热衷的追求中。\par

\SourceRecord{Translated by Peter Holmes and Robert Ernest Wallis, and revised by Benjamin B. Warfield.；Nicene and Post-Nicene Fathers, First Series；Vol. 5.；Edited by Philip Schaff.；Buffalo, NY: Christian Literature Publishing Co.,；1887.；<http://www.newadvent.org/fathers/15071.htm>.}

% structure: p0001 source=h1 target=section reason=page-title

\section{论婚姻与贪欲（卷二）}

% structure: p0002 source=h2 target=subsection reason=relative-heading-rank; base=h2

\subsection{预备说明}

% structure: p0003 source=h3 target=subsubsection reason=relative-heading-rank; base=h2

\subsubsection{（一）摘自\Person{奥古斯丁}的\WorkTitle{未完成的反儒利安著作}序言}

我写了一篇题为\WorkTitle{论婚姻与贪欲}的论文，呈献给\Person{瓦莱里乌斯}伯爵，因为我听说贝拉基派异端者指责我们谴责婚姻。在那篇论文中，我尽可能严谨而准确地指出了婚姻的善与肉欲之恶的区别——这种恶被婚姻的贞洁所善用。当那位杰出人士收到我的论文后，他通过一份简短的文件寄给我一些摘录，这些摘录选自贝拉基派异端者\Person{儒利安}的一部著作。在这部著作中，他认为合适将自己的回答扩展为四卷，以回应我先前提到的仅为一卷的论文\WorkTitle{论婚姻与贪欲}。我不知道是谁为我们提供了这些摘录：他显然有意将选择限于\Person{儒利安}著作的第一卷。应\Person{瓦莱里乌斯}的要求，我立即撰写了回答这些摘录的文字。于是，我以同一标题写了第二卷；\Person{儒利安}对此以八卷的文字予以回答，这超出了他饶舌的能力。\par

% structure: p0005 source=h3 target=subsubsection reason=relative-heading-rank; base=h2

\subsubsection{（二）摘自\Person{奥古斯丁}致\Person{克劳狄乌斯}书函[第CCVII号]}

\Quote{凡读过我这第二卷书的人（该卷同第一卷一样，呈献给\Person{瓦莱里乌斯}伯爵，并实际上两者都是为他而撰写的），都会发现有些地方我并未回答\Person{儒利安}，而是更多针对那位从\Person{儒利安}的书中摘录的人，并且他没有按摘录时的顺序编排。他认为有必要对其编排做相当大的改变，很可能是为了借此方法将明显属于他人的思想据为己有。}\par

% structure: p0007 source=h2 target=subsection reason=relative-heading-rank; base=h2

\subsection{本书正文}

\Person{奥古斯丁}在此后一卷中驳斥了某些句子，这些句子是由某位不知名作者从\Person{儒利安}为反对他前一部著作\WorkTitle{论婚姻与贪欲}而出版的四卷书中的第一卷摘录的；这些句子应\Person{瓦莱里乌斯}伯爵的请求被转呈给他。他维护了天主教关于原罪的教义，驳斥了对手的吹毛求疵和诡辩，并特别指出该教义与臭名昭著的摩尼教异端是多么不同。\par

% structure: p0009 source=h2 target=subsection reason=relative-heading-rank; base=h2

\subsection{第一章——引言}

我最亲爱且受尊崇的儿子\Person{瓦莱里乌斯}，我无法告诉你，当听说你在军务及你当之无愧的高位所赋予的职责，以及繁忙的政治生活中，仍对反对异端者的天主圣言的见证怀有如此热切而真挚的兴趣时，我心中涌起何等大的喜乐。读完\Person{瓦莱里乌斯}阁下您的信函（您在信中承认了我献给您的那卷书），我感到有必要再写这一卷；因为您要我注意我的兄弟兼同主教\Person{阿利比乌斯}受命向我陈述的事，即异端者就我书中的若干段落所引发的争论。我不仅从我兄弟的叙述中得知此事，也读了他提交的摘录，这些摘录是您在他离开\Place{拉文纳}后亲自寄往\Place{罗马}的。发现我们对手的夸夸其谈后，我很容易在这些摘录中看到，我定意，靠主的助佑，用我能掌握的全部真理和圣经权威来回应他们的嘲讽。\par

% structure: p0011 source=h2 target=subsection reason=relative-heading-rank; base=h2

\subsection{第二章——在此章及随后四章中，他提出了需审议的曲解摘录}

我现在回答的这份文件以此标题开头：\Quote{\Emph{摘自\Person{奥古斯丁}所著书的标题，我据此从书中摘录了若干段落}}。我由此看出，将这份书面文件寄给\Person{瓦莱里乌斯}阁下您的人，想从他不具名的书中摘录，依我判断，是为了更快得到答复，以免耽搁您的急切。在考虑他指的是哪些书之后，我想必定是\Person{儒利安}在寄往\Place{罗马}的书信中提到的那些，该信的一份副本同时送到了我手中。因为他在那里说：\Quote{他们甚至断言，现在争论中的婚姻不是天主建立的——这一说法可在\Person{奥古斯丁}的一部著作中读到，我最近以一篇四卷的论文作了回答。}我相信，这些就是摘录所出自的书。那么，也许更好的做法是，我有意驳斥并推翻他扩展为四卷的全部著作。但我极不愿意延迟我的答复，因为您自己也毫不拖延地寄来了需要我回答的书面陈述。\par

% structure: p0013 source=h2 target=subsection reason=relative-heading-rank; base=h2

\subsection{第三章——同上继续}

他引用了并试图反驳的、出自我的书中的那几段话，就是下面这些；那本书我已送给你，你非常熟悉。\Quote{他们因过度仇恨我们，不断地断言我们谴责婚姻以及天主藉男女结合创造人的神圣程序，因为我们主张从这种结合出生的人染有原罪，并且我们不否认，无论其父母是谁，他们若不在基督内重生，就仍在魔鬼的权势之下。}  现在，他在引用我的这些话时，故意省略了我所引用的宗徒的见证，因为他感到那见证的重大意义对他压力太大。因为在说了人生来就染有原罪之后，我立刻引用了宗徒的话：\Quote{罪恶藉一人进入了世界，死亡藉罪恶而进入；这样死亡就传给了众人，因为众人都犯了罪。}\BibleRef{罗 5:12} 正如我已提到的，他省略了宗徒的这段话，然后结束了我刚才引用的其他评论。因为他太清楚，对于所有忠信的天主教徒的心灵和良心来说，我采用而他却省略的宗徒的这些话是多么可接受——这些话是如此直接、如此清晰，以致这些标新立异的异端者尽其黑暗曲折的曲解所能，竭力模糊和败坏它们的力量。\par

% structure: p0015 source=h2 target=subsection reason=relative-heading-rank; base=h2

\subsection{第四章 同上。}

但他又加上了我的另一些话，我说过：\Quote{他们也不反省：婚姻的善并不因源自婚姻的原恶而受到更多指责，正如奸淫和私通的恶也不能因出自它们的自然善而得到原谅。因为罪恶无论来自前者还是后者，都是魔鬼的工作；同样，人无论由前者还是后者所生，都是天主的工作。} 在这里，他又遗漏了一些词句，他害怕天主教徒的耳朵。因为要说到此处引用的这些话，我们先前曾说过：\Quote{因此，因为我们肯定这一教义，它包含在古老不变的天主教的信仰规则中，这些提出新奇而乖戾教义的人，否认婴儿有任何需要在重生水泉中被洗净的罪，因他们的不信或无知而毁谤我们，仿佛我们谴责婚姻，仿佛我们断言那本是天主的工作、即由婚姻所生的人，是魔鬼的工作。} 所有这些段落他都略过了，只引用了随后如上所述的话。现在，在他省略的词句中，他害怕那个适用于天主教会所有人心、诉诸那自古以坚定不移的声音被坚定确立并传承下来的信仰的句子，并且最强烈地激起他们对我们的敌意。这个句子是：\Quote{他们否认婴儿有任何需要在重生水泉中被洗净的罪。} 因为所有人带着婴儿跑向教会，世界上没有别的原因，只是为了通过第二次出生的重生，洗净他们第一次自然出生时所染上的原罪。\par

% structure: p0017 source=h2 target=subsection reason=relative-heading-rank; base=h2

\subsection{第五章 同上。}

然后他又回到我们先前引用的话：\Quote{我们主张，从这种结合出生的人染有原罪；我们不否认，无论其父母是谁，他们若不在基督内重生，就仍在魔鬼的权势之下。} 为什么他又要再次引用我们的话，我说不出；他刚才已经引用过了。然后他继续引用我们关于基督的话：\Quote{祂不愿从两性的同一结合而生。} 但在这里，他又悄悄地忽略了我放在这些话之前的话；我的整个句子是：\Quote{好藉祂的恩宠，他们能从黑暗的权势中被救出，并被移入那不愿从两性的同一结合而生的祂的国度。} 请注意，我恳求你，他回避了哪些话，表现出一种完全反对那藉我们的\Quote{主耶稣基督}而来的天主恩宠的态度。他十分清楚，将婴儿排除在宗徒的话所带给他们的利益之外，是极度的不正直和不虔敬；宗徒论及天主圣父时说：\Quote{祂拯救了我们脱离了黑暗的权势，把我们移入了祂爱子的国度。}\BibleRef{哥 1:13} 毫无疑问，这就是他宁愿省略而不引用这些话的原因。\par

% structure: p0019 source=h2 target=subsection reason=relative-heading-rank; base=h2

\subsection{第六章 同上。}

他随后引用了我们的一段话，其中我们说了：\Quote{因为如果没有先前犯罪，就不会有这种引发羞耻的私欲偏情——无耻之徒竟无耻地赞扬它；至于婚姻，即使无人犯罪，它仍然会存在：因为子女的生育本可以在没有这种病患的情况下完成。} 至此他引用了我的话；但他却不敢引用紧接着的话——\Quote{在那个贞洁生命的身体内，虽然实际上没有它，但在\InnerQuote{这个死亡的躯体}内就无法完成。} 他不愿意完成我的句子，而是将其截断，因为他害怕宗徒的呼喊，我的话语提醒了他：\BibleQuote{我这个人真不幸呀！谁能救我脱离这个死亡的躯体呢？天主的恩宠，藉着我们的主耶稣基督。} \BibleRef{罗 7:24} 因为这个死亡的躯体在犯罪前的乐园里并不存在；因此我们曾说：\Quote{在那个贞洁生命的身体内}，即乐园的生命，\Quote{子女的生育本可以在没有这种病患的情况下完成，而现在在这个死亡的躯体中，没有它就无法完成。} 然而，宗徒在到达我们刚才引用的关于人的悲惨和天主的恩宠之处以前，首先说过：\BibleQuote{我发觉在我的肢体中另有一条法律，与我理性的法律交战，并把我俘虏到那在我肢体中的罪恶的法律。} 然后他喊道：\BibleQuote{我这个人真不幸呀！谁能救我脱离这个死亡的躯体呢？天主的恩宠，藉着我们的主耶稣基督。} 在这个死亡的躯体中——它在犯罪前的乐园里当然不是如此——肯定没有\BibleQuote{在我们的肢体中另有一条法律，与我们理性的法律交战}——这条法律如今，即使在我们不愿意、不同意，且不使用我们的肢体去满足它欲望的时候，仍然居住在这些肢体中，扰乱我们抵抗和反抗的理性。而这种冲突本身，虽然不涉及定罪，因为它没有完成犯罪，但仍然是\Quote{不幸的}，因为它没有平安。我想，我已经足够清楚地表明，此人在引用我的话时有特殊的目的和方法：他以反驳的方式引用了它们，有时通过删除介于中间的文字来打断我句子的连贯性，有时通过去掉结尾词语来缩短它们；而我认为我已经充分解释了他这样做的原因。\par

% structure: p0021 source=h2 target=subsection reason=relative-heading-rank; base=h2

\subsection{第七章 [III.] 奥古斯丁引用来朱利安前言中的一段话。（参见\Quote{\WorkTitle{未完成之作}，卷一，73}。）}

\Quote{我们当代的教师们，至圣的兄弟，他们是那个可耻派别的煽动者，该派别如今因狂热而发热；他们决心不惜以整个教会的毁灭为代价，来图谋伤害和诋毁那些以其神圣热忱点燃了他们的人；他们丝毫没有考虑到，他们赐予那些人何等大的荣誉，因为他们表明，那些人的名声只能与公教一同被毁灭。因为，若有人说人内有自由意志，或说天主是出生者的创造者，他立刻被斥为赛勒斯提乌斯和白拉奇派。为了避免被称为异端者，他们变成摩尼教徒；如此，在逃避一种虚假的耻辱的同时，他们招致真正的谴责；就像动物一样，人们在狩猎中用染色的羽毛包围它们，以惊吓并驱赶它们进入网中；可怜的畜生没有理性天赋，因此它们被一种虚惊一起推入真正的毁灭。}\par

% structure: p0023 source=h2 target=subsection reason=relative-heading-rank; base=h2

\subsection{第八章。—— 奥古斯丁反驳上引用段落。}

好了，现在，无论你是谁，说了这一切，你所说的绝非真实；绝非，我重复；你大受欺骗，或者你意在欺骗他人。我们不否认自由意志；但是，正如真理所宣称的：\Quote{如果圣子使你们自由，你们就的确是自由的了。}\BibleRef{若 8:36}正是你们自己恶意否认这位解放者，因为你们在俘虏状态中把一种虚妄的自由归于自己。你们是俘虏；因为\Quote{人被谁制胜，就是谁的奴隶；}\BibleRef{伯后 2:19}除了藉着伟大解放者的恩宠，没有人能从这种束缚的锁链中得释放，任何活着的人都不能免于此锁链。因为\Quote{藉着一个人，罪进入了世界，藉着罪，死亡；这样死亡就传给了众人，因为众人都犯了罪。}\BibleRef{罗 5:12}因此，天主是那些如此出生者的创造者，以致所有没有藉着再生而拥有那唯一解放者的人，都从那一人进入定罪。因为祂被描述为\Quote{陶匠，从同一团泥中，制造一个器皿为尊贵，出于祂的仁慈，另一个为卑贱，出于审判。}\BibleRef{罗 9:21}因此教会咏唱着\Quote{仁慈和审判。}所以，当你们说：\Quote{如果有人声称，要么人中有自由意志，要么天主是出生者的创造者，他立即被定为塞勒斯提乌斯派和白拉奇派，}你们只是在误导自己和他人；因为公教信仰是说这些的。然而，若有人说人中有自由意志可以无需祂的帮助而正当地敬拜天主，并且有人说天主是出生者的创造者，以致否认婴儿需要有人将他们从魔鬼的权势中救赎出来：那人才是被定为塞勒斯提乌斯和白拉奇的追随者。因此，人里面有自由意志，以及天主是出生者的创造者，这两点我们都承认。你们并非仅仅因为说了这些就成了塞勒斯提乌斯派和白拉奇派。但你们实际说的是，任何人都有足够的意志自由，无需天主帮助而行善，并且婴儿没有经历那种\Quote{从黑暗的权势中被救出来，被迁到天主的国度里}\BibleRef{哥 1:13}的改变；因为你们这样说，所以你们是塞勒斯提乌斯派和白拉奇派。那么，你们为什么躲在一个共同教条的遮蔽之下进行欺骗，隐藏使你们获得派别之名的特殊罪行？并且，为了用一个骇人的术语恐吓无知者，你们论到我们说：\Quote{为了避免被称为异端者，他们变成了摩尼教派？}\par

% structure: p0025 source=h2 target=subsection reason=relative-heading-rank; base=h2

\subsection{第九章——公教徒坚持原罪的教义，因此远非摩尼教派}

那么，请稍听片刻，看看这个问题涉及什么。公教徒说，人性是由良善的天主作为创造者所创造的，原是良善的；但因罪而败坏，需要医生基督。摩尼教派肯定人性不是由天主创造为良善的，也没有因罪败坏；而是人是由永远存在的两种本性——一种良善，一种邪恶——的混合，由永恒黑暗的王子所形成的。白拉奇派和塞勒斯提乌斯派说，人性是由良善的天主创造为良善的；但在婴儿出生时仍如此健全和健康，以致他们在那个年龄不需要基督的药物。那么，请在你的教义中认出你的名字；不要再将属于别人的名字和教义强加给驳斥你们的公教徒。因为真理拒绝双方——摩尼教派和你们。对摩尼教派，它说：\Quote{你们没有念过：那起初造人的，是造男造女，并且说：\InnerQuote{因此，人要离开父母，与妻子连合，二人成为一体}吗？这样，他们不是两个，而是一体。所以，天主所结合的，人不可分开。}\BibleRef{玛 19:4}在此，基督表明天主既是人的创造者，也是夫妻婚姻的结合者；而摩尼教派否认这两点。然而，对你们，祂说：\Quote{人子来，是为寻找并拯救那迷失了的。}\BibleRef{路 19:10}但你们，作为卓越的基督徒，回答基督：\Quote{如果祢来是为寻找并拯救那迷失了的，那么祢就不是为婴儿来的；因为他们并没有迷失，而是生来就在得救状态中；去找年长的人吧；我们给你一条出自祢话语的规则：\InnerQuote{健康的人不需要医生，而是有病的人才需要。}\BibleRef{玛 9:12}}恰好，摩尼教派说人的本性中混合了邪恶，他至少希望他的善良灵魂能被基督拯救；而你们却主张婴儿中没有需要被基督拯救的东西，因为他们已经是安全的。因此，摩尼教派以其可憎的指责围攻人性，而你们以其残忍的赞美围攻人性。因为凡相信你们赞美的人，绝不会将他们的婴孩带到救主面前。怀着如此不敬的观点，你们无畏地面对那为你们设下的、为引起有益的恐惧并像对待人（而非你们那只被彩色羽毛包围、驱入猎网的可怜动物）一样对待你们的法令，这有什么用呢？你们本该持守真理，并且因对真理的热诚而无所畏惧；但事实上，你们逃避恐惧的方式是，如果你们害怕，你们更愿意逃离恶者的网罗，而不是投入其中。你们的公教母亲使你们恐惧，是因为她为你们以及别人因你们而担心；如果她藉着在国中拥有权威的儿子的帮助，使你们畏惧，她这样做不是出于残忍，而是出于爱。然而，你们是非常勇敢的人；你们认为害怕人是懦夫的行为。那么，敬畏天主吧；不要如此顽固地试图颠覆公教信仰的古老根基。尽管我甚至希望你们那种激烈的性情在目前的情况下会对人的权威有一点点恐惧。我希望，我说，它宁可在胆怯中战栗，也不要在胆大妄为中灭亡。\par

% structure: p0027 source=h2 target=subsection reason=relative-heading-rank; base=h2

\subsection{第10章 [IV.]——如何驳斥对手的诡辩。}

现在我们来看他在其选集中所拼接的其余部分。但我应当采取什么步骤呢？我是否应当列出他的每一段话并加以回答，或者，省略一切公教信仰所包含的、在我们之间没有争议的内容，只处理并驳斥那些他偏离真理正途、并试图将白拉奇异端的毒枝嫁接到公教树干上的陈述？这无疑是较为容易的做法。但我认为，我不可忽略一种可能的情况，就是任何人在读了本书之后，若没有通读他提出的所有论据，可能会认为我不愿提出他所依据的段落，而藉这些段落，我所反驳的虚假陈述才真正得以推论出来。因此，我希望读者能善意地留意并考量这部小作品中出现的两类贡献——也就是说，他所提出的全部内容，以及我对他所作的回答。\par

% structure: p0029 source=h2 target=subsection reason=relative-heading-rank; base=h2

\subsection{第11章——魔鬼是罪的作者，而非本性的作者。}

那位将上述文件呈送给您的爱德的人，在其序言中以这样的标题开头：\Quote{反对那些谴责婚姻、并将其果实归于魔鬼的人。}因此，这并非反对我们，我们既不谴责婚姻——我们甚至以其正当的次序恰当地称赞它——也不将其果实归于魔鬼。因为婚姻的果实是从中正常生育的人，而非伴随他们出生的罪。人之所以在魔鬼的权势之下，并非因为他们是人（在这方面他们是婚姻的果实），而是因为他们是罪人（这涉及罪的传递）。因为魔鬼是罪的作者，而非本性的作者。\par

% structure: p0031 source=h2 target=subsection reason=relative-heading-rank; base=h2

\subsection{第12章——厄娃的名字意为生命，并且是教会的一个伟大圣事。}

现在请看他在那段文字中的其余部分，他自以为在其中找到了与上述标题相符的内容，损害我们的利益。他说：\Quote{天主，祂用泥土塑造了亚当，又用他的肋骨造了厄娃，\BibleRef{创 2:22}并说：\InnerQuote{她要称为生命，因为她是一切活人的母亲。}}其实，圣经并非如此记载。但这与我们何干？因为我们的记忆常常在字句上出错，而意思仍然保持。而且，给厄娃命名的是她的丈夫，而非天主，这名字应表示\Emph{生命}；因为经上如此写道：\BibleQuote{亚当给他妻子起名生命，因为她是一切活人的母亲。}但他很可能理解为圣经藉亚当（作为他的先知）证明天主给了厄娃这个名字。因为称她为生命、一切活人的母亲，这其中包含教会的一个伟大圣事，要详细谈论它会耽搁太久，且对我们当前的任务并无必要。宗徒所说的\BibleQuote{这是一个伟大圣事：但我是指基督和教会而言}，也正是亚当在说\BibleQuote{为此，人要离开父亲和母亲，并要依附于他的妻子，二人成为一体}时所说的话。然而，主耶稣在福音中提及天主曾就厄娃说了这话；这无疑是因为天主藉着人宣布了那人作为预言所说的话。现在请看这摘录文件中接下来的内容：他说：\Quote{藉着这原始的名字，祂表明了女人被造是为了什么工作；祂相应地说了：\BibleQuote{你们要生育繁殖，充满大地。}\BibleRef{创 1:28}}在我们当中，有谁否认女人是由主天主、一切美善的恩惠创造者为生育的工作而预备的呢？再看他又说了什么：\Quote{因此，\BibleQuote{天主造了他们一男一女}\BibleRef{创 1:27}，给他们配备了适于生育的肢体，并命令身体由身体产生；并且祂保证他们有能力完成这工作，以祂在创造中所使用的德能造出一切存在之物。}我们也承认这是公教信理，同样地，他立即附加的那段也是如此：\Quote{因此，如果后代只能通过性别而来，而性别只能通过身体而来，身体则通过天主而来，谁能犹豫而不承认生育力应归功于天主呢？}\par

% structure: p0033 source=h2 target=subsection reason=relative-heading-rank; base=h2

\subsection{第13章——白拉奇论点，旨在证明魔鬼对婚姻的果实没有任何权利。}

在这些真实而公教的陈述之后——它们确实包含在圣经中，尽管他并非以公教的精神、怀着公教的心志来引证它们——他立即向我们介绍白拉奇和塞勒斯提乌斯的异端，他写下上述评论正是为了这个目的。请仔细留意下面的话：\Quote{现在你们这些人说：\InnerQuote{我们不否认，无论由何父母所生，他们仍在魔鬼的权势之下，除非他们在基督内重生}，请向我们指出，在性别中魔鬼能认作属于自己的是什么，从而他可以（用你们的话说）正当地声称他所产生的果实为自己的财产。是性别的\Emph{差异}吗？这在于天主所造的身体之内。是它们的\Emph{结合}吗？这结合在原始祝福的特恩中与在设立中同样得到认可。因为天主的声音说：\BibleQuote{男人要离开他的父亲和母亲，并要依附于他的妻子，二人成为一体}\BibleRef{创 2:24}又是天主的声音说：\BibleQuote{你们要生育繁殖，充满大地}\BibleRef{创 1:28}或者，或许是它们的\Emph{生育力}吗？但这正是婚姻被设立的原因。}\par

% structure: p0035 source=h2 target=subsection reason=relative-heading-rank; base=h2

\subsection{第14章 [V.]——唯独私欲偏情（在婚姻中）不是来自天主的。}

你看到了他向我們提出的問題：魔鬼能在兩性中找到甚麼可稱為自己的，以至於使那些無論是由何种父母所生的人，除非在基督內重生，否則就落在他的權下？他又問我們，是否我們將兩性間的差異、或是他們的結合、或是他們的生育能力歸於魔鬼？我們回答說，這些性質沒有一樣是屬於魔鬼的，因為兩性的差異屬於父母的\Quote{器官}；兩者的結合屬於生育子女；他們的生育能力屬於婚姻制度的祝福。但這一切都是出於天主；然而他不願提及那\Quote{肉體的私慾，不是出於父，而是出於世界}；\BibleRef{若一 2:16}而魔鬼被稱為\Quote{這世界的首領}\BibleRef{若 14:30}主身上沒有肉體的私慾偏情，因為主不是藉著它而成為人來到人間。因此，祂親自說：\Quote{這世界的首領來了，在我內一無所有}\BibleRef{若 14:30}——一無所有，即沒有罪；既沒有從出生而來的，也沒有在生命中加添的。在他所提及的所有生育的自然善中，他，我重複說，不願意提及這私慾偏情的具體事實，連婚姻也為之羞愧，而婚姻卻以所有以上所述的善為榮。為甚麼父母特別的工作甚至對自己的子女也隱藏起來？除非因為他們不可能沒有可恥的私慾而從事可讚美的生育。正因為此，那些最初遮蓋自己裸體的人也感到羞愧。\BibleRef{創 3:7}他們身體的這些部分從前並不羞恥，反而配得作為天主的化工而受讚美。他們在感到羞愧時穿上了遮蓋；而當他們在違抗造物主之後，感到自己的肢體不服從自己時，他們感到了羞愧。我們這位摘錄者也對這私慾偏情感到了羞愧。因為他提到了兩性的差異；他也提到了他們的結合，以及他們的生育能力；但這最後伴隨私慾的事，他羞於提及。難怪那些空談者對父母本身也羞於思考的事感到羞恥，而父母對自己的職責卻是如此關心。\par

% structure: p0037 source=h2 target=subsection reason=relative-heading-rank; base=h2

\subsection{第十五章——人因罪而生在魔鬼的權下；當亞當犯罪時，我們眾人都在他內。}

然後他繼續問：\Quote{那麼，為甚麼天主所創造的人卻在魔鬼的權下呢？}他似乎是從我的一句話中找到自己問題的答案。\Quote{因為罪，}他說，\Quote{不是因為本性。}然後他根據我的回答來構建他的答案：\Quote{但是正如沒有兩性就沒有後代，同樣沒有意志就沒有罪。}是的，的確如此。因為\BibleQuote{正如罪藉著一人進入了世界，死亡藉著罪也進入了世界；這樣死亡就傳到了眾人，因為眾人都在他內犯了罪。}\BibleRef{羅 5:12}因那一個人的邪惡意志，眾人在他內都犯了罪，因為眾人都是那一個人，因此他們各自承受了原罪。\Quote{因為你斷言，}他說，\Quote{他們在魔鬼權下的原因是因為他們是由兩性結合而生的。}我明確宣稱，他們是因違命而在魔鬼權下，而且他們對這違命的參與是由於他們由所說的兩性結合而生的，而這種結合甚至無法在沒有可恥的私慾偏情的情況下完成其可敬的職能。事實上，這也是至可紀念的\Place{米蘭}教會主教\Person{安博羅削}所說過了的，他在解釋為何基督在肉體中的誕生完全沒有罪惡的過錯時，給出的理由是祂的成孕不是兩性結合的結果；而在這種結合中成孕的人類中，沒有一個是無罪的。他的原話是：\Quote{因此，祂作為人，受了各種試探的考驗，並按人的樣式承受了這一切；然而，由於祂是由聖神所生，祂遠離了一切罪惡。因為每個人都是說謊者，除了天主以外，沒有人是無罪的。因此，}他補充說，\Quote{這一直是常被注意到的，即顯然沒有一個由男人和女人所生的人，也就是說，藉著他們身體的結合，是無罪的；因為誰若無罪，誰也免於這種成孕。}那麼，\Person{佩拉糾}和\Person{塞萊斯提}的弟子們，你們敢稱這個人為摩尼教徒嗎？就像異端者\Person{若維尼安}所做的那樣，當這位聖主教維護真福瑪利亞在生育後仍保持永久的童貞以對抗此人的不敬時。然而，如果你們不敢稱他為摩尼教徒，那麼當我們以同樣的理由和同樣的見解維護公教信仰時，你們為甚麼稱我們為摩尼教徒呢？但如果你們要嘲弄那位最忠信的人在這個問題上持有摩尼教的謬誤，那就沒有辦法了，你們只能盡情享受你們的嘲弄，從而更加完全地充滿若維尼安的尺度；至於我們，我們可以與這樣一位天主之人一同耐心忍受你們的嘲弄和譏諷。然而，你們的異端首領\Person{佩拉糾}極力讚揚\Person{安博羅削}的信德和對聖經知識的極度純淨，甚至聲稱連敵人也無法挑剔他。那麼，看看你們已經走到了何地步，並停止跟隨\Person{若維尼安}大膽的腳步。然而，那個人雖然過分讚揚婚姻，將其與神聖的童貞等同，卻從未否認基督的必要性來拯救那些從婚姻中出生、甚至剛出母腹的嬰兒，並將他們從魔鬼的權下贖回。但是你們否認這點；因為我們反對你們，以維護那些尚不能為自己說話的嬰兒以及公教信仰的基礎本身，你們就嘲弄我們是摩尼教徒。但現在讓我們看看接下來的話。\par

% structure: p0039 source=h2 target=subsection reason=relative-heading-rank; base=h2

\subsection{第十六章——魔鬼是我們罪的作者，而不是我們本性的作者。}

他向我们提出另一个问题，说：\Quote{那么，你承认谁是婴儿的作者？真天主吗？}我回答说：\Quote{是的，真天主。}然后他评论说：\Quote{但祂没有创造邪恶；}又问：\Quote{我们是否承认魔鬼是婴儿的创造者？}然后他再次回答说：\Quote{但\Emph{他}没有创造人性。}然后他好像以这个推论结束主题：\Quote{既然结合是邪恶的，我们身体的状态是堕落的，因此你把我们的身体归于一个邪恶的创造者。}我对此的答复是，我没有把我们的身体归于邪恶的创造者，而是把我们的罪归于他；由于这些罪，结果在我们身体里，就是说，在天主所造的一切中，原本都是可尊敬且令人喜悦的，但在男女的交往中却产生了羞耻之事，以至于他们的结合不是那种在无玷生命中可能有的结合，而是我们在羞愧中在这个死亡的身体里所看到的结合。\Quote{但是天主，}他说，\Quote{在性别上分开了祂要在运作中结合的东西。所以身体结合来自祂，身体的创造也首先来自祂。}我们已经对这个陈述提供了答复，当我们说这些身体来自天主。但至于这些身体肢体的不顺服，这是来自于肉体的贪欲，这贪欲\Quote{不是源自圣父的。}\BibleRef{若一 2:16}他接着说，认为\Quote{不可能从如此多美好的事物，如身体、性别和它们的结合中，结出邪恶的果实；或者，人类被天主造出来，目的是按你所坚持的，以合法权利被魔鬼占有。}现在我们已经肯定，他们不是由于他们是人而被占有，这个名称属于他们的本性，魔鬼不是其作者；而是由于他们是罪人，这个名称是本性的那个缺陷的结果，魔鬼是那缺陷的作者。\par

% structure: p0041 source=h2 target=subsection reason=relative-heading-rank; base=h2

\subsection{第十七章 [VII.]——佩拉基乌斯派不羞于赞美贪欲，尽管他们羞于提及它的名称。}

但在如此多美好的事物名称中，如身体、性别和结合，他从未提到肉体的贪欲。他保持沉默，因为羞耻；然而以一种奇特的羞耻之耻（如果可以这么说的话），他不羞于赞美他羞于提及的东西。现在观察他如何更喜欢通过迂回而非直接提及来指向他的目标。\Quote{在那之后，那人，}他说，\Quote{因自然的欲求认识了他的妻子。}再看，他拒绝说\Quote{因肉体的贪欲认识了他的妻子}；而是用了\Quote{因自然的欲求}这个短语，这让我们可以理解那神圣而可敬的意愿，即愿意生育子女，而不是那种连他都如此羞于启齿的贪欲，以至于他宁愿对我们使用模棱两可的语言，也不愿用明确的词语表达他的思想。现在，他所说的\Quote{因自然的欲求}是什么意思？难道得救的愿望和生养、抚育子女的愿望不也是自然的欲求吗？并且不也是理性的，而非贪欲的？然而，既然我们能确定他的意图，我们相当肯定他意在通过这些词语表示生殖器官的贪欲。难道你认为这些词语不就是无花果叶，其遮盖下的不是别的，正是他感到羞耻的东西吗？因为正如他们古时用叶子缝制起来\BibleRef{创 3:7}作为遮盖的腰带一样，这人编织了一张迂回之网来隐藏他的意思。让他编织他的陈述：\Quote{但当那人因自然的欲求认识了他的妻子时，圣经说：厄娃怀了孕，生了一个儿子，给他起名叫加音。但是，}他补充说，\Quote{亚当说了什么？让我们听：我从天主得到了一个人。}\BibleRef{创 4:1}因此，很明显他是天主的工程，圣经见证了他从天主领受的事实。谁能怀疑这一点？谁能否认这个陈述，尤其如果他是一个天主教基督徒？人是天主的作品；但肉体的贪欲（如果没有罪先存在，人本会通过生殖器官被孕育，这些器官与其他肢体一样顺从于安静正常的意志）不是源自圣父的，而是属于世界的。\BibleRef{若一 2:16}\par

% structure: p0043 source=h2 target=subsection reason=relative-heading-rank; base=h2

\subsection{第十八章。——同上续。}

但现在，我求你，更仔细地看，观察他如何设法找到一个名字来再次掩盖他耻于揭示的东西。\Quote{因为，}他说，\Quote{亚当藉他肢体的力量生了那个孩子，而不是由于功绩的差异。}现在我承认，我不明白他后面那句\Emph{不是由于功绩的差异}是什么意思；但当他说\Quote{藉他肢体的力量}时，我相信他想要表达他耻于公开明确说出的话。他更喜欢用\Quote{藉他肢体的力量}这个短语，而不是说\Quote{藉肉体的贪欲。}显然——即使他没有想到这一点——他暗示了某样与该主题明显相关的东西。因为当一个人的肢体不适当顺从人的意志时，还有什么比这些肢体更有力量？即使它们被节德或贞操所约束，它们的使用和控制也不在任何人的权力之下。因此，亚当藉我们作者所谓的\Quote{他肢体的力量}生了儿子，在他生他们之前，他在犯罪后曾为此羞愧。然而，如果他从未犯罪，他就不会藉肢体的力量，而是藉肢体的服从来生他们。因为他自己将会有力量来统治它们，使它们服从他的意志，如果他本人也以同样的意志只服从于一位更有大能者作为臣属。\par

% structure: p0045 source=h2 target=subsection reason=relative-heading-rank; base=h2

\subsection{第十九章 [VIII.]——佩拉基乌斯派误解圣经中的\Quote{种子}。}

他接着说：\Quote{过了一会儿，圣经又说：\BibleQuote{亚当认识了他的妻子厄娃，她生了一个儿子，给他起名叫舍特，说：\InnerQuote{主给了我另一个后代代替亚伯尔，因为加音杀了他。}}}然后他补充道：\Quote{据说神性本身生出了这个后代，以此证明性结合是祂的安排。}这个人不明白圣经所记载的；因为他以为圣经说\BibleQuote{主给了我另一个后代代替亚伯尔}的原因无他，就是使人以为天主激发了亚当对性交的欲望，借此生出种子倾入女人子宫。他完全不知道圣经所说的并非\BibleQuote{主给了我后代}（如他所用的意义），而只是\BibleQuote{主给了我一个儿子。}确实，亚当不是在他性交后排出种子时说这话，而是在他妻子分娩后，因天主的恩赐得到了儿子时说的。因为仅仅是排出种子作为性结合的最后快乐，若不伴随婚姻的真正果实——怀孕和生育，有什么可取之处呢（也许对好色之徒和那些如宗徒所禁止的人而言：\BibleQuote{在邪淫的欲望中占有了自己的器皿}\BibleRef{得前 4:5}）？\par

% structure: p0047 source=h2 target=subsection reason=relative-heading-rank; base=h2

\subsection{第二十章——原罪源自人类种子的败坏状态。}

然而，我这样说，绝非暗示我们要在至高真实的天主之外另寻造物主，无论是人类种子还是由种子而来的人本身；而是说，如果没有罪在先，种子本应借着肢体对意志命令的安静、正常的服从而从人身发出。现在我们面前的问题不是人类种子的本性，而是它的败坏。\Emph{本性}以天主为作者；原罪源自它的败坏。如果种子本身没有败坏，那么《智慧篇》中的这句话是什么意思：\BibleQuote{他们本来是邪恶的种子，生来的恶意不能改变，因为他们的种子从起初就是受诅咒的}\BibleRef{智 12:10}？无论这些话特别应用于谁，它们都是关于人类的。那么，为什么每个人的恶意是生来的，他的种子从起初就受诅咒，难道不是因为\BibleQuote{藉着一个人，罪进入了世界，藉着罪，死亡也进入了世界，这样死亡就传给了众人，因为众人都犯了罪}\BibleRef{罗 5:12}吗？但哪里有人\BibleQuote{邪恶的意图永远不能改变}，除非因为他自己不能做到，只能靠天主的恩宠？没有这帮助，人除了像伯多禄宗徒所描述的\BibleQuote{生来就是畜牲，好像生来该受捉拿和毁灭}\BibleRef{伯后 2:12}，还能是什么呢？因此，保禄宗徒在某一处，同时考虑到两种状况——我们出生时所带的天主义怒，以及我们得救的恩宠——说：\BibleQuote{我们从前也都在他们中间，放纵肉体的私欲，随从肉体和心思的意愿行事，生来就是义怒之子，和别人一样。然而富于慈悲的天主，因着他爱我们的大爱，竟在我们因过犯死了的时候，使我们同基督一起活过来——因恩宠你们得救了。}\BibleRef{弗 2:3}那么，什么是人的\BibleQuote{生来的恶意}和\BibleQuote{从起初受诅咒的种子}？什么是\BibleQuote{生来就是畜牲，好像生来该受捉拿和毁灭}？什么是\BibleQuote{生来就是义怒之子}？这是在亚当所受造的本性的状况吗？绝对不！因为他的纯洁本性在他里面被败坏了，便通过自然遗传在所有人中持续这种状态，并且仍然在持续；以至于除非藉着天主在我们的主耶稣基督内的恩宠，没有人能从这败坏中得救。\par

% structure: p0049 source=h2 target=subsection reason=relative-heading-rank; base=h2

\subsection{第二十一章 [IX.]——使人多产的是良善的天主，败坏果实的是魔鬼。}

因此，这个人在下一段话中是什么意思？他说到\Person{诺厄}和他的儿子们：\Quote{他们蒙受了祝福，正如\Person{亚当}和\Person{厄娃}一样；因为天主对他们说：\BibleQuote{你们要生育繁殖，充满大地，治理大地。}}\BibleRef{创 9:1}对全能者的这些话，他补充了自己的话，说：\Quote{现在，你们称之为魔鬼的快乐，在上述的婚姻夫妇中也被采用；它继续存在于其制度的良善及其附带的祝福中。因为毫无疑问，以下的话是针对\Person{诺厄}和他的儿子们说的，关于他们与妻子的身体结合——这结合此时已因习惯而不可改变地固定了：\BibleQuote{你们要生育繁殖，充满大地。}}我们无需多言重复先前的论证。这里的问题是，我们本性中的败坏，其良善已被扭曲，而魔鬼正是这败坏的作者。本性本身的良善——其作者是天主——不是我们要考虑的问题。天主从未从败坏和扭曲的本性中撤回祂的怜悯和良善，以致剥夺人的生育力、生机、健康，以及身心实体本身，还有感官、理性，以及食物、营养和成长。此外，祂\BibleQuote{使太阳升起，光照恶人，也光照善人；降雨给义人，也给不义的人}\BibleRef{玛 5:45}；人类本性中一切良善都来自良善的天主，即使是在那些不会从恶中得到拯救的人身上也是如此。\par

% structure: p0051 source=h2 target=subsection reason=relative-heading-rank; base=h2

\subsection{第二十二章——我们将为我们所做的感到羞耻，还是为天主所做的？}

然而，此人（指对方）在他的段落中谈论的是\Emph{快乐}，因为快乐甚至可以是可敬的；至于产生羞耻的\OriginalTerm{肉身的贪情}{carnal concupiscence}或\OriginalTerm{情欲}{lust}，他只字未提。但在随后的话语中，他暴露了自己对羞耻的易感性；他无法掩饰本性强力地规定的事。他说：\Quote{还有那句话：\BibleQuote{因此人要离开自己的父母，依附自己的妻子，二人成为一体。}} 然后，在引用了天主的这些话之后，他继续发表自己的见解，说：\Quote{先知为了在行为中表达信德，几乎危及了端庄。} 这是何等的供认！多么清晰，被真理的力量从他身上逼出！先知，为了表达信德于行为，在说\BibleQuote{二人成为一体}时，几乎危及了端庄，他这话是指男女的性结合。为什么要提出理由，说明先知在表达天主的作为时，竟如此接近危及端庄呢？难道人的作为不应产生羞耻，而必须被夸耀，而天主的作为却必须产生羞耻吗？难道在陈述和表达天主的作为时，先知的爱或辛劳得不到尊荣，反而他的端庄受到危及吗？那么，天主做了什么，竟使他的先知羞于描述？而且，一个更重大的问题是：一个人能否对任何不是人、而是天主在人内造成的作为感到羞耻？然而，工人在所有情况下都以他们所能的劳作和勤勉，努力避免对自己亲手所做的工感到羞耻。但事实是，我们正是对那件使原始人感到羞耻而遮掩腰部的同一件事感到羞耻。那是罪的惩罚；那是罪的瘟疫和烙印；那是罪的诱惑和真正的燃料；那是\OriginalTerm{在我们肢体中的法律}{law in our members}，与我们的心神的法律交战；那是来自我们自己、反对我们自己的反叛，通过我们叛逆的肢体，由最公义的报复归还给我们。正是这个使我们感到羞耻，而且公义地感到羞耻。若非如此，我们若因我们肢体的混乱而羞愧，不是因我们自己的过犯或惩罚，而是因天主的作为，那还有什么比这更忘恩、更不虔敬的呢？\par

% structure: p0053 source=h2 target=subsection reason=relative-heading-rank; base=h2

\subsection{第二十三章 [X.]——伯拉纠派主张天主在\Person{亚巴郎}和\Person{撒辣}的情况下，激起贪情作为来自天上的恩赐。}

关于\Person{亚巴郎}和\Person{撒辣}，他（指对方）也说了许多，虽然毫无意义，谈到他们如何按照恩许得了一个儿子；最后他提到了\Emph{\OriginalTerm{贪情}{concupiscence}}这个词。但他没有加上通常的短语\Quote{属于肉身的}，因为正是这引起了羞耻。然而，有时因贪情也可夸耀，因为有一种心神对肉身的贪情\BibleRef{迦 5:17}，也有一种智慧的贪情。因此，他说：\Quote{如今你们确实将这为生育所不可或缺的贪情界定为天生邪恶；那么，它如何会在老年人中被从天赐予而被激起呢？那么，如果你们能够，就请澄清，你们视为由天主作为恩赐所赐予的，竟属于魔鬼的作为。} 他这样说，就好像肉身的贪情从前在他们身上不存在，好像是天主将其赐予了他们。无疑，它（贪情）本已存在于这死亡的身体中；然而，生育能力是缺乏的，而天主是生育的作者；这实际上是当主乐意赐予恩赐时给予的。然而，我们绝不要确认他以为我们想要说的，即\Person{依撒格}是在没有性结合的激情下受生的。\par

% structure: p0055 source=h2 target=subsection reason=relative-heading-rank; base=h2

\subsection{第二十四章 [XI.]——新生儿违反天主的哪种盟约。割损的价值是什么。}

但请告诉我们，若不在第八天受割损，他的灵魂为何会从人民中被剪除？若没有原罪在其中，他如何能犯下如此大的罪，如何如此得罪天主，以致因别人对他疏忽而受如此严厉的惩罚？关于婴孩割损，天主的诫命如此说：\BibleQuote{凡未受割损的男子，即在第八天未割损包皮的，他的灵魂应从人民中被剪除，因为他违背了我的盟约。}\BibleRef{创 17:14}请告诉我们，若他能够，那孩子如何违背了天主的盟约——一个八天大的无辜婴孩，就他本人而言？然而天主或圣经绝无谎言。事实是，他所违背的天主盟约并非这条命令割损的，而是那条禁止吃树的；那时\BibleQuote{因一人，罪进入了世界，藉着罪，死亡也进入了世界；于是死亡传给了众人，因为众人都犯了罪。}\BibleRef{罗 5:12}在他身上，第八天的割损即藉着要降生成人的中保的圣事，象征了这种赎罪。因为正是藉着对那将在肉身来临、为我们而死、并在第三天（即安息日后，为第八天）复活的基督的同一信仰，古圣们也得救了。因为\BibleQuote{祂被交出来，是为我们的过犯；祂复活了，是为我们成义。}\BibleRef{罗 4:25}自从割损在神的子民中设立以来（当时那是信德之义的标记），它甚至对婴孩也有效，用以象征原初罪过的洗净，正如圣洗从设立之时起同样开始用于更新人。并非在割损之前没有因信德成义；因为\Person{亚巴郎}自己还在未受割损时，就因信德而成义，成为那些也要效法他信德的万民之父。\BibleRef{罗 4:10}然而，在古时，因信德成义的圣事奥秘隐藏在各种方式中。但仍是那同一的对中保的信德拯救了古圣们——无论大小——不是那\BibleQuote{生子为奴}的旧约，\BibleRef{迦 4:24}不是那并未被赐予以能赐予生命的法律，\BibleRef{迦 3:21}而是藉着我们的主耶稣基督的天主恩宠。\BibleRef{罗 7:25}因为我们相信基督已降生成人，他们相信祂将要来临；同样，我们相信祂已死了，他们相信祂将要死；我们相信祂已从死者中复活，他们相信祂将要复活；而我们和他们一样相信，祂将来要审判活人死人。因此，让此人不要藉着为人的功德作拙劣辩护，而为人的得救设置障碍，因为我们都是在罪中出生，又由那唯一无罪出生者被救赎。\par

% structure: p0057 source=h2 target=subsection reason=relative-heading-rank; base=h2

\subsection{第二十五章 [XII.]——\Person{奥古斯丁}并非原罪的发明者}

\Quote{这身体的性结合，}他说，\Quote{连同那炽热、连同那快感、连同精液的射出，是由天主所造，就其本身而言是可称赞的，因此应当认可；它甚至有时成为对虔诚之人的一份大恩赐。}他明确而分别地重复了\Quote{连同那炽热}、\Quote{连同那快感}、\Quote{连同精液的射出}。然而他没有敢于说\Quote{连同情欲}。这是为什么，难道不是因为他羞于命名那他不羞于称赞的吗？确实，对虔诚之人而言，子女的繁盛是一项恩赐；但并非那使肢体羞愧的激动——这激动，若本性健康，我们便不会感觉到，虽然败坏的本性如今经历它。正因如此，由此而生的人需要重生，才能成为基督的肢体；而生他的人，即使已经重生，也需要从那因罪的律法存在于这死亡之身中的东西中被释放。既然情况如此，他如何能继续说：\Quote{因此你必须不得不承认你所杜撰的原罪已被消除}？并非我杜撰了原罪——公教信仰自古持有这道理；而你否认它，无疑是个标新立异的异端。在天主的审判中，除非他们在基督内重生，否则所有人都在魔鬼的权下，生于罪中。\par

% structure: p0059 source=h2 target=subsection reason=relative-heading-rank; base=h2

\subsection{第二十六章 [XIII.]——孩子绝非由欲望形成}

但当他论到\Person{亚巴郎}和\Person{撒辣}时，他接着说：\Quote{如果你们确实肯定他们身上自然之欲很强烈，却没有后裔，我的答复是：造物主所应许的，造物主也赐予了；所生的孩子不是同寝的工程，而是天主的工程。的确，那位用尘土造了第一个人的，也由种子塑造所有人。因此，如同作为材料的尘土本身不是人的作者；同样，那形成并混合精液元素的性欲能力，并不完成造人的全部过程，而是从自然宝库中向天主呈现材料，天主屈尊用以造人。}现在他整个这段话，除了他说精液元素是由性欲形成和混合的那部分以外，如果他愿意诚恳地用它来辩护公教含义，本可正确表述。然而，对于我们这些完全知道他试图从中得出什么结论的人来说，他确实以扭曲的方式正确地说着。这个对普遍真理的例外陈述——我不否认属于这一段——之所以不真实，是因为所涉及的血肉私欲偏情的快感并不形成精液元素。这些元素已经存在于身体中，并由创造身体本身的同一真天主所形成。它们并非从色欲快感获得存在，而是与之一起被激发和排出。确实，这种快感是否伴随着两性精液元素在子宫内的混合，或许妇女能从其内心感受来判断；但我们不应将无聊的好奇心推得太远。然而，那种我们必须为之羞愧的私欲偏情，其羞耻给了我们隐秘肢体以可耻的名称\OriginalTerm{阴部}{pudenda}，在罪恶进入之前，身体在乐园生活期间并不存在；而是在罪恶之后，它开始存在于\InnerQuote{这死身体}内，肢体反叛以报复人自身的不服从。没有这种私欲偏情，婚姻对偶的生育功能完全可能实现：正如许多劳累的工作由我们其他肢体的顺从运作来完成，没有任何淫欲的灼热；因为它们只是由意志的指示所驱动，而非被私欲偏情的激情所煽动。\par

% structure: p0061 source=h2 target=subsection reason=relative-heading-rank; base=h2

\subsection{第二十七章——白拉奇派争论说天主有时在愤怒中关闭子宫，在被平息时打开它。}

请仔细考虑他其余的言论：他也说，\Quote{这也}由宗徒的权威所证实。因为当圣保禄论及死者复活时，他说道：\Quote{愚昧人哪，你所播种的，若不死去，就不能活过来。}\BibleRef{格前 15:36}此后，\InnerQuote{但天主随自己的意愿给它一个身体，并给每个种子它自己的身体。}因此，天主，他说，\Quote{已经给人类种子，如同给其他一切，它自己固有的身体，没有一个明智或虔诚之人会否认。那么，你如何证明任何人生来有罪呢？我求你，请反思，这自然之罪的论断被何等绞索扼杀。但是，来吧，}他说，\Quote{请你对自己更温和些，我求你。相信我，连你也是天主造的；然而必须承认，一个严重的错误已感染了你。因为还能提出什么更亵渎的意见，不是天主没有造人，就是祂把人造给了魔鬼；或者，至少是魔鬼塑造了天主的肖像，即人——这显然不仅是荒谬，而且是亵渎？那么，}他说，\Quote{天主是如此资源贫乏，如此缺乏任何体统，以至于除了魔鬼欺骗他们后可能注入使他们败坏的东西，就没有任何东西可以赐予圣人们作为赏报吗？然而，你想知道，即使是那些并非圣人的，也能证明天主赐予了他们生育子女的能力吗？当\Person{亚巴郎}在外邦人中因恐惧而说他的妻子\Person{撒辣}是他的妹妹时，据说当地国王\Person{阿彼默肋客}把她带走，欲与她过夜。但天主，那圣女的荣誉由祂保守，在梦中显现给\Person{阿彼默肋客}，制止了王室的放肆；威胁他若侵犯妻子必死。然后\Person{阿彼默肋客}说：\InnerQuote{主啊，你要杀害一个无辜和正义的民族吗？他们不是告诉我他们是兄妹吗？}因此\Person{阿彼默肋客}清早起来，拿了一千块银子，以及羊、牛、仆婢，给了\Person{亚巴郎}，并将他的妻子安然送回。但\Person{亚巴郎}为\Person{阿彼默肋客}向天主祈祷；天主就治好了\Person{阿彼默肋客}和他的妻子及他的婢女。}现在他为何如此冗长地叙述这一切，你可以在他补充的这几句话中找到：\Quote{天主，}他说，\Quote{因\Person{亚巴郎}的祈祷，恢复了她们的生育能力，这能力甚至从最卑微的婢女的子宫中被夺走；因为天主关闭了\Person{阿彼默肋客}家中每一个子宫。\BibleRef{创 20:18} 现在，}他说，\Quote{思考一下，这应否被称为一种自然之恶，有时天主在愤怒时取走，平息时恢复。祂，}他说，\Quote{既造虔诚者的子女，也造不虔者的子女，因为他们是父母这一事实属于那以天主为造物主而喜乐的本性，而他们的不虔属于他们欲望的败坏，这作为自由意志的后果临到每一个人。}\par

% structure: p0063 source=h2 target=subsection reason=relative-heading-rank; base=h2

\subsection{第二十八章 [十四]——奥斯定对此论点的答复。其与圣经的关系。}

现在，我们必须回答他那冗长的陈述：在他从圣经引用的段落中，丝毫没有提及那种可耻的情欲；我们说过，在我们原祖的幸福状态中，他们赤身露体并不羞惭时，这种情欲并不存在于他们的身体内。\BibleRef{创 2:25} 他从宗徒那里引用的第一段经文是关于谷物的种子，这些种子必须先死去才能复活。不知为何，他不愿完整引述这节经文。他从中引用的全部内容是：\BibleQuote{愚钝的人哪，你所撒的并不复活；} 但宗徒补充说：\BibleQuote{除非死了。} \BibleRef{格前 15:36} 然而，据我判断，这位作者希望那些既不理解也不记得圣经的人，将这段只论及谷物种子的经文理解为人的种子。事实上，他不仅删减了这句话，省略了\Quote{除非死了}这一从句，而且还省略了以下词语，宗徒在这些词语中解释了他所指的是何种种子；因为宗徒补充说：\BibleQuote{并且你所撒的，并不是那将来要有的形体，而是赤裸的谷粒，或许是麦子，或许是别的谷粒。} \BibleRef{格前 15:37} 他省略了这些，并用宗徒随后所写的文字来结束他的上下文：\BibleQuote{但天主随自己的心意给它一个形体，并给各种子各有的形体。} 就好像宗徒在谈论人类的同居时说了\BibleQuote{愚钝的人哪，你所撒的并不复活}，以便我们理解人的种子是由天主复活的，而不是由人在同居中生儿育女。因为他先前曾说：\Quote{性快感并不能完成造人的全部过程，而只是从自然的宝藏中向天主提供材料，天主愿意用这些材料来造人。} 然后他加上引文，好像宗徒肯定如下：\Quote{愚钝的人哪，你所撒的并不复活——也就是说，不是靠你自己复活；而是天主从你的种子中造人。} 就好像宗徒没有说过那些中间的话，而这位作者故意跳过；又好像宗徒的目的就是如此谈论人的种子：\Quote{愚钝的人哪，你所撒的并不复活；但天主随自己的心意给种子一个形体，并给各种子各有的形体。} 确实，在宗徒的话之后，他引进了自己的评论，大意如此：\Quote{因此，如果天主给人种子，如同给其他一切种子一样，它自己适当的形体——这是任何有智慧或敬虔的人都不会否认的——} 完全好像宗徒在所讨论的段落中谈论的是人的种子。\par

% structure: p0065 source=h2 target=subsection reason=relative-heading-rank; base=h2

\subsection{第二十九章——继续同一论题。\Person{奥斯定}亦坚持天主在出生时造人。}

尽管我特别留意这一点，但我未能发现他通过这种欺骗性地使用圣经能得到什么帮助，只想把宗徒作为证人，并借由他证明，我们也主张的，即天主从人的种子造人。所以，由于没有直接找到合适的经文，他欺骗性地篡改了这段经文；无疑是因为担心，如果宗徒在这段经文中似乎说到谷物的种子，而非人的种子，我们就会立即通过他的这种行为而知道如何驳斥他：不是把他当作一个洁身自好、主张克制意志的人，而是厚颜无耻、宣扬放纵淫逸的人。但是，实际上，他从农夫撒在田里的种子本身就可以被驳倒。因为我们为何不能设想，天主可能赐予在乐园幸福状态中的人，与祂赐予谷物的种子相同的对待其自身种子的方式，即前者可以在没有任何可耻情欲的情况下播下，生育的肢体单纯服从意志的倾向，正如后者可以在没有任何可耻情欲的情况下撒下，农夫的手只是按他的意志而动作呢？当然，有一个区别，即父母生育子女的欲望比农夫装满谷仓的欲望更高尚。再者，为什么全能的造物主，凭其无玷的遍在，以及从人的种子中任意造物的能力，不能以同样的方式在女人身上运作，正如祂现在所造的，按照祂的旨意像在土地中处理谷物种子那样，使得有福的母亲在没有情欲的情况下怀孕，在没有产痛的情况下生产，因为（在那种幸福状态中，在尚未成为这死亡的身体、而是那生命的身体中）女人在接受种子时没有任何可羞耻的事，正如生产时没有任何痛苦的事一样？凡拒绝相信这一点，或不愿设想，在人未有罪恶之前生活在乐园的幸福状态中，我们所描述的这种状态不可能蒙天主的美意和仁慈所允许的人，必须被视为可耻快感的爱好者，而非可欲繁殖力的赞美者。\par

% structure: p0067 source=h2 target=subsection reason=relative-heading-rank; base=h2

\subsection{第三十章[第十五]——\Person{阿彼默肋客}及其家室案之考察。}

再者，关于他从灵感的历史中引用的有关阿彼默肋客的一段，以及天主选择封闭他全家每一个子宫，使妇女不能生育，后来又开启子宫使她们能生育的事，这与论点有何关系？这与我们现在所争论的羞耻的情欲有何关系？难道天主剥夺了这些妇女的这种感觉，又在祂喜欢的时候还给她们吗？然而，惩罚是她们不能生育，祝福是她们能够生育，依照这必朽血肉的方式。因为天主不会赐予这死亡的身体这样的祝福，只有乐园中那个生命之体在罪恶进入之前才可能拥有；也就是说，没有情欲的刺激而怀孕，没有剧烈的疼痛而生育。但我们为何不这样想：既然圣经确实说每一个子宫都被封闭，这件事伴随着某种疼痛，以致妇女无法忍受同房，而天主在祂的义怒中施加了这种疼痛，并在祂的仁慈中移除了它？因为如果情欲要被移除作为生育的障碍，它本应从男人而不是从女人身上移除。因为即使激发她的情欲消失，女人也可以凭她的意志履行她在同房中的职责，只要男人身上不缺少那激发他的情欲；除非，也许（正如圣经告诉我们，连阿彼默肋客自己也被治愈了），他会说男性的情欲被归还给了他。然而，如果他真的失去了这个，有什么必要让天主警告他不要接近亚巴郎的妻子呢？事实是，阿彼默肋客被认为是治愈了，因为他的家人从击打他们的灾祸中痊愈了。\par

% structure: p0069 source=h2 target=subsection reason=relative-heading-rank; base=h2

\subsection{第三十一章 [XVI.]——为何天主继续创造人类，祂知道他们将生于罪中。}

现在让我们来看他提出的三个句子，他说，再也没有比这三个更亵渎的了：\Quote{要么天主没有造人，要么祂为了魔鬼造人；或者，魔鬼塑造了天主的肖像，也就是人。}现在，第一句和最后一句，即使是他自己，如果他不鲁莽和乖僻，也必须承认我们从未说过。争论只限于他放在中间的那一句。关于这一点，他错得如此之远，竟以为我们说过天主为了魔鬼造人；好像在天主从人类父母所创造的人身上，祂的关怀、目的和安排是，借由祂的工作，魔鬼应拥有那些他无法自己制造的人作奴隶。天主禁止任何敬虔的信仰，即使是最幼稚的，竟怀有这种情感！按照祂自己的良善，天主造了人——第一个人没有罪，所有其他的人都在罪之下——为了祂深奥思想的目的。因为正如祂完全知道如何处理魔鬼本身的恶意，而祂所做的，无论魔鬼多么不义和邪恶，都是公义和良善的；也正如祂不因预知他会邪恶而不愿创造他；同样，关于整个人类，虽然没有人不带着罪的污染而生，但至善的天主总是在行善，使一些人成为所谓的\Quote{仁慈的器皿}，恩宠将他们与那些是\Quote{义怒的器皿}的人分别出来；而使另一些人成为所谓的\Quote{义怒的器皿}，为要彰显祂的荣耀的丰富在仁慈的器皿上。\BibleRef{罗 9:23}那么，让这个反对者去和宗徒辩论吧，我用的是他的话；不，是去和那位陶工辩论，宗徒禁止我们反驳祂，正如著名的经文所说：\BibleQuote{你这个人哪，你是谁，竟敢向天主抗辩？受造之物岂能对造它的说：你为什么这样造了我？难道陶工没有权柄从同一团泥中，造一个器皿为尊贵，另一个器皿为卑贱吗？}\BibleRef{罗 9:20}那么，这个人要争辩说，义怒的器皿不在魔鬼的管辖之下吗？或者，因为他们在这管辖之下，他们是由另一位造物主所造，而不是那位造仁慈器皿的？或者祂是用别的材料、而不是从同一团泥中造的？那么，他可以反对说：\Quote{因此天主造这些器皿是为了魔鬼。}好像天主不知道如何利用这些器皿，甚至魔鬼本身，来促进祂自己良善公义的工作。\par

% structure: p0071 source=h2 target=subsection reason=relative-heading-rank; base=h2

\subsection{第三十二章 [XVII.]——天主并非祂所创造之人中邪恶的作者。}

那么，天主喂养灭亡之子，祂左边的山羊，\BibleRef{玛 25:33}是为了魔鬼，养育和衣着他们是为了魔鬼，\BibleQuote{因为祂使太阳上升，光照恶人，也光照善人；降雨给义人，也给不义的人}\BibleRef{玛 5:45}？那么，祂创造恶人，正如祂喂养和养育恶人一样；因为祂藉创造所赐予他们的，属于本性的良善；而祂藉饮食和营养所给予他们的成长，自然是祂仁慈的帮助，不是给予他们的邪恶，而是给予那本性的良善，祂在祂的良善中创造了这本性。因为就他们是人而言——这是那个本性的善，其作者和创造者是天主；但就他们生于罪中，除非重生否则注定灭亡而言，他们属于那从起初受咒诅的种子，\BibleRef{智 12:11}由于原始背命的过错。然而，这位造者连义怒的器皿也利用这过错为善，为要彰显祂的荣耀的丰富在仁慈的器皿上；\BibleRef{罗 9:33}并且没有人可以将他因恩宠而得救归功于任何自己的功绩，因为他属于同一团泥；而是\BibleQuote{夸耀的，应当因主而夸耀。}\BibleRef{格后 10:17}\par

% structure: p0073 source=h2 target=subsection reason=relative-heading-rank; base=h2

\subsection{第三十三章 [XVIII.]——虽然天主造了我们，但除非祂在基督里重造我们，我们都会灭亡。}

从此宗徒的、大公的信德最真实、最稳固的原则出发，我们面前的这位作者与白拉奇派一同背离了。他不愿承认人生来处于魔鬼的统治之下，以免婴孩被带到基督那里，从黑暗权势中被解救出来，并迁入祂的王国。\BibleRef{哥 1:13} 如此，他成了遍及普世的教会的控告者；在这教会中，到处都有婴孩在受洗时首先被驱魔，别无他故，只因要让这世界的首领从他们中被赶出去。\BibleRef{若 12:31} 因为他们必然被魔鬼所占有，作为震怒的器皿，既然他们从亚当而生，除非他们在基督内重生，并藉恩宠被迁入怜悯的器皿的王国。然而，他在攻击这一最稳固的真理时，试图避免公然攻击整个基督教会。因此，他只对我一人提出上诉，并以责难和劝诫的口吻说：\Quote{但天主也造了你，尽管必须承认一种严重的错误已感染了你。} 好吧，我感谢地承认天主也造了我；然而，如果祂只从亚当造了我，而没有在基督内再造我，我必与震怒的器皿一同灭亡。但是，此人既然被白拉奇的异端所附，他并不相信这一点；如果他确实坚持如此重大的错误直到最后，那么，不是他，而是天主教徒，将能看清那错误的大小和性质——这错误没有简单地感染他，而是完全毁灭了他。\par

% structure: p0075 source=h2 target=subsection reason=relative-heading-rank; base=h2

\subsection{第三十四章——白拉奇派主张正当使用同居是好的，由此所生亦是好的。}

我请求你现在注意以下的话。他说：\Quote{然而，在婚姻中怀上的孩子本性上是好的，我们可以从宗徒的话学到，当他谈到那些人，\BibleQuote{他们离开了女人的自然用途，在他们的淫欲中焚烧，男与男行可耻之事。}\BibleRef{罗 1:27} 在这里}，他说，\Quote{宗徒显示女人的用途既是自然的，也是在其方式上可称赞的；滥用则在于行使自己的意志，违反制度的正派使用。} 所以，他说，\Quote{在那些正当使用它的人中，私欲偏情在其种类和方式上得到称赞；而它的过度，即放纵之人所沉溺的，受到惩罚。} 事实上，正当天主以火罚惩罚索多玛人的滥用时，祂增强了亚巴郎和撒辣因年迈而变得无力的生育能力。因此，他继续说，\Quote{你认为必须因生育器官的力量而指责它们，因为索多玛人借此沉溺于罪中，那么你也必须谴责像饼和酒这样的受造物，因为圣经告诉我们，他们也滥用了这些恩赐。因为主藉祂的先知厄则克耳的口说：\InnerQuote{\BibleQuote{这些，而且是你的姊妹索多玛的罪；她与她的女儿骄傲，因食物丰足和安逸享乐而自满，却没有扶助穷人和急需者的手。}} \BibleRef{则 16:49} 因此，选择吧}，他说，\Quote{你宁愿要哪一个：要么将人类身体的性结合归咎于天主的作为，要么认为像饼和酒这样的受造物同样邪恶。但如果你偏向这后者，你就证明自己是\OriginalTerm{摩尼教徒}{Manichean}。然而，真理是：他若在自然的私欲偏情中遵守节制，就是善用善物；他若不遵守节制，就是滥用善物。那么，你的声明是什么意思呢？}他问道，\Quote{当你说：\InnerQuote{婚姻的善不再因由此而来的原罪而受指控，正如奸淫和淫乱的恶不能因由此而生的自然之善而被原谅一样。}用这些话}，他说，\Quote{你承认了你所否认的，又使你承认的无效；你除了要让人不能理解之外，别无所图。请给我指出没有性结合的身体上的婚姻。否则就给这行为赋予一个名称，并将夫妻结合指定为善或恶。你无疑会回答，你已经将婚姻定义为善。那么，若婚姻是善——若人是婚姻的善果；若这果实，作为天主的工程，不能是恶，由善的媒介从善而生——那原初的恶在哪里，它已被如此多的先前承认所搁置？}\par

% structure: p0077 source=h2 target=subsection reason=relative-heading-rank; base=h2

\subsection{第三十五章——他答复朱利安努斯的论点。什么是女人的自然使用？什么是非自然使用？}

我对此挑战的答复是：不仅婚姻所生的儿女，连淫乱所生的儿女，就其是天主创造的工程而言，都是好的；但就原罪而言，他们都是在第一个亚当之下被判罪而生，不仅那些生于淫乱的人，同样那些生于婚姻的人，除非他们在第二个亚当——即基督——内重生。至于宗徒论及恶人时所说：\BibleQuote{男人放弃了女人的自然之用，彼此欲火中烧，男与男行可耻之事；}\BibleRef{罗 1:27} 他并非指婚姻之用，而是指\Quote{自然之用}，希望我们明白，藉着为此目的而造的肢体，两性如何能够结合以生育。由此可见，即使人与娼妓结合而使用这些肢体，这种使用仍是自然的。然而，它不可称赞，反而有罪。但就身体任何非为生育目的的部分而言，即使人用其与自己的妻子结合，也是违反本性的可耻行为。事实上，同一宗徒此前\BibleRef{罗 9:26}关于妇女曾说：\BibleQuote{连他们的女人也把自然的之用变为违反本性的；}然后关于男人他补充说，他们因放弃女人的自然之用而行可耻之事。因此，这里所说的\Quote{自然之用}并不是为了赞美婚姻结合，而是用以指那些比男人使用女人（即使不合法，仍是自然的）更为不洁和罪恶的可耻行为。\par

% structure: p0079 source=h2 target=subsection reason=relative-heading-rank; base=h2

\subsection{第三十六章【XXI.】——天主造的自然原是好的：救主修复被败坏的自然。}

现在我们并不因有些人奢侈酗酒就指责饼和酒，也不因贪婪和吝啬就非难金子。同样，我们也不因身体那引起羞耻的情欲而斥责夫妻之间可敬的结合。因为前者（即无羞耻的结合）在犯任何罪之前是完全可能的，夫妻结合时不致脸红；而后者（即情欲）是在犯罪之后产生的，他们因极度羞愧不得不将其隐藏。\BibleRef{创 3:7} 因此，自那以后，在所有夫妻结合中，无论他们多么正当合法地使用这种恶，都永远有必要避免被人看见这种行为，从而承认那必然引起羞耻的事——好的事当然不会使人羞愧。这样，我们不知不觉地注意到两个不同的事实：那值得称赞的两性结合以生育子女的好，与那羞耻的情欲的恶，因这恶，后代必须重生才能免于定罪。因此，一个人若虽带着那引起羞耻的情欲却合法同居，便是以善用恶；若非法同居，则是以恶用恶——因为那使善恶之人皆脸红的事，确切地说，是恶而非善。我们更应相信那曾说\BibleQuote{我知道，在我内，即在我肉身内，没有善}\BibleRef{罗 7:18}的人，而非称那使自己如此承认其为恶的东西为善的人；但若他不觉羞愧，便是加增无耻的更恶。因此，我们正确地宣称：\Quote{婚姻之善并不因由之而来的原罪而受指责，正如奸淫与淫乱之恶不能因由之而生的自然之善而得开脱；}因为无论由婚姻或淫乱所生的人性，都是天主的工程。如果这本性是恶的，就不该被生；如果没有恶，就不必重生。并且（为将二者归于同一判断），如果人性是恶的，它就不该被拯救；如果它内没有任何恶，它就不必被拯救。因此，主张本性非善的人，是说造物主不善；而主张本性内无恶的人，是在其败坏的状态中剥夺了它一位仁慈的救主。由此可知，在人的出生中，既不能因善的造物主所形成的善而原谅淫乱，也不能因仁慈的救主需要医治的恶而指责婚姻。\par

% structure: p0081 source=h2 target=subsection reason=relative-heading-rank; base=h2

\subsection{第三十七章【XXII.】——若没有同居就没有婚姻，那么也没有不羞耻的同居。}

\Quote{请你指出，}他说，\Quote{任何没有性结合的肉身婚姻。}我指不出任何没有性结合的肉身婚姻；但同样，他也指不出任何没有羞耻的性结合。然而在乐园中，若没有先犯罪，固然会有无两性结合就不生育，但这种结合必然没有羞耻；因为在性结合中，肢体将安静地顺从，而非那产生羞耻的肉身情欲。所以婚姻是好的，在其中人藉有序的受孕而生；婚姻的果实也是好的，即如此而生的人本身；但罪是恶的，人人带着这恶而生。然而，天主造人并仍造人；但\BibleQuote{藉一人，罪进入了世界，藉罪，死亡；这样死亡就传给了众人，因为众人都犯了罪。}\BibleRef{罗 5:12}\par

% structure: p0083 source=h2 target=subsection reason=relative-heading-rank; base=h2

\subsection{第三十八章【XXIII.】——约维尼安从前称天主教徒为摩尼教徒；亚略派也称天主教徒为撒伯里乌派。}

\Quote{藉着你们这种新的辩论方式，}他说，\Quote{你们既宣称自己是公教徒，又庇护摩尼教，因为你们既称婚姻是极大的善，又称它是极大的恶。}他完全不知道自己在说什么，或是假装不知道。或者他不理解我们说的话，或者他不希望人们理解。如果他不理解，那是他被错误的先入之见所阻碍；如果他不希望人们理解我们的意思，那么他就是在固执地维护自己的错误。几年前试图创立一个新异端的\Person{约维尼安}也曾宣称，公教徒庇护摩尼教徒，因为他们反对他，将神圣的童贞置于婚姻之上。但此人肯定会回答说，他不赞同约维尼安对婚姻和童贞的漠不关心。我并不是说这就是他们的观点；但这些新异端者必须承认，约维尼安对公教徒使用摩尼教徒的指责这一事实表明，这种做法并不新鲜。我们宣称婚姻是善，不是恶。但正如亚略派指控我们是撒伯里乌派，尽管我们并不像撒伯里乌派那样说圣父、圣子和圣神是同一的，而是像公教徒那样肯定圣父、圣子和圣神有同一的本性；同样，佩拉奇派也以摩尼教徒来指责我们，尽管我们并不像摩尼教徒那样宣称婚姻是恶，而是像公教徒那样肯定邪恶临到第一对男女，即第一对夫妻，并从他们传到所有人。然而，正如亚略派在避开撒伯里乌派时，落入了更坏的伙伴，因为他们胆敢分割的不是天主圣三的位格，而是本性；同样，佩拉奇派在努力逃离摩尼教徒的致命错误时，走向了另一个极端，被证明对婚姻的果实持有比摩尼教徒更恶劣的观点，因为他们相信婴儿不需要基督作为他们的医生。\par

% structure: p0085 source=h2 target=subsection reason=relative-heading-rank; base=h2

\subsection{第三十九章［第二十四章］——无论出自何种父母，人都是生而有罪且能得救赎的。}

接着他说：\Quote{你们断定，若生于奸淫，人便无罪；若生于婚姻，人便非无罪。因此，你们的断言等于说，从奸淫关系中可能产生自然之善，而从婚姻中却实际传递了原罪。}他现在试图将正确的话曲解，但在明眼读者面前是徒劳的。我们绝不说，生于奸淫的人无罪。但我们断言，一个人，无论生于婚姻还是奸淫，在某些方面是好的，因为这是由于天主，万物的创造者；然而我们补充说，他因原罪而承袭了某种恶。因此，我们的说法——\Quote{从奸淫父母也可能产生自然之善，但从婚姻也传递了原罪}——并不等于他试图归咎于我们的：即生于奸淫的人无罪，生于婚姻的人非无罪；而是说，无论生于何种状况，人都因原罪而有罪；而无论哪种状况的后代，都可以因自然之善而通过重生获得自由。\par

% structure: p0087 source=h2 target=subsection reason=relative-heading-rank; base=h2

\subsection{第四十章［第二十五章］——\Person{奥古斯丁}拒绝接受他提出的两难困境。}

\Quote{这两个命题，}他说，\Quote{一个为真，一个为假。}我的回答和这一指控一样简短：两者都是真的，没有一个是假的。\Quote{真实的是，}他接着说，\Quote{奸淫之罪不能因所生之人而得到原谅；因为奸淫者所犯的罪属于意志的败坏；但他们所生的后代则有利于对生育的赞美。如果有人偷了麦种去播种，长出的庄稼并不因此而变坏。}\Quote{当然，}他说，\Quote{我谴责偷窃者，但我赞美麦粒。同样，我宣告那从种子慷慨的果实而生的人是无罪的；正如宗徒所说：\InnerQuote{天主随自己的意思给它一个身体，并给每个种子一个它自己的形体。}\BibleRef{格前 15:38}但同时，我谴责那以不正当方式使用神圣规定而犯下奸淫之罪的恶人。}\par

% structure: p0089 source=h2 target=subsection reason=relative-heading-rank; base=h2

\subsection{第四十一章［第二十六章］——佩拉奇派争辩说，如果婚姻是善，原罪就不能来自婚姻。}

在这之后，他接着说了以下的话：\Quote{如果确实从婚姻中产生了恶，那么婚姻可以被指责，不，不能被原谅；并且你们将婚姻的工作和果实置于魔鬼的权下，因为凡是恶的原因本身就没有善。然而，从婚姻所生的人，其起源不在于婚姻的责难，而在于其精原；这些精原的原因在于身体的状态；任何人若滥用这些身体，便是损害其善的功绩，而非其本性。因此，}他辩称，\Quote{显然，善不是恶的原因。所以，}他继续说，\Quote{如果原恶甚至从婚姻而来，那么恶的原因就是婚姻的结合；那么那恶果所由来的东西本身也必然是恶的；正如主在福音中所说：\InnerQuote{凭果实可认出树来。}\BibleRef{玛 7:16}那么，}他问道，\Quote{当你们说婚姻是善，却又宣称从它而来的只有恶时，你们如何认为自己值得注意？显然，婚姻是有罪的，因为原罪是从它们推演出来的；而且它们也是不可辩护的，除非它们的果实被证明是无辜的。但它们是受到辩护的，并被宣告为善；因此它们的果实被证明是无辜的。}\par

% structure: p0091 source=h2 target=subsection reason=relative-heading-rank; base=h2

\subsection{第四十二章——佩拉奇派试图借赞美天主的化工来消除原罪；婚姻就其本质和制度而言，不是罪的原因。}

我对这一切已有答复；但在给出答复之前，我希望读者注意，这些人的意见所导致的结果是，他们认为婴儿不需要救主，因为他们认为婴儿完全没有需要被拯救的罪。这种对真理的巨大歪曲，如此敌视天主的伟大恩宠——这恩宠是藉着我们的主耶稣基督赐予的，祂\BibleQuote{来寻找并拯救失丧的人}\BibleRef{路 19:10}——试图通过赞美天主的工程来潜入缺乏智慧者的心灵；也就是说，通过赞美人性、人的种子、婚姻、性交、婚姻的果实——这一切都是美好的事物。我不会说他还赞美了情欲；因为他甚至羞于提及它，所以他所赞美的似乎是别的东西，而不是\Emph{它}。他用这种方法，不区分加诸于本性的邪恶与本性自身的良善；他确实没有表明本性是健全的（因为这不真实），但他却不允许本性的疾病状态得到医治。因此，我方的第一个主张，即由奸淫所生的好东西，即人本身，并不能为奸淫关系的罪开脱，他承认这是真的；对于这一点，我们之间没有争论，他甚至（恰当地）用窃贼偷种种子而长出极好庄稼的比喻来辩护和强化它。然而，我方的另一个主张，即\Quote{婚姻之善不能为源自婚姻的原罪受责备}，他却不承认为真；如果他同意这一点，他就不会是佩拉纠异端者，而是公教基督徒。他说：\Quote{当然，}他说，\Quote{如果邪恶来自婚姻，它就可以受责备，甚至不能被原谅；而你把它的工作和果实置于魔鬼的权势之下，因为凡邪恶之因本身都没有善。}除此之外，他还设计了其他论证，表明善不可能成为恶的原因；由此他推断，婚姻是善，不是恶的原因；因此，没有人能从婚姻中出生而有罪的状态并需要救主；就好像我们说婚姻是罪的原因，尽管事实上在婚姻中出生的人并非没有罪而生。婚姻的设立不是为了犯罪，而是为了生育子女。因此，主对婚姻状态的祝福如此说：\BibleQuote{你们要生育繁殖，充满大地。}\BibleRef{创 1:28}然而，从婚姻传递给子女的罪不属于婚姻，而是属于加诸于人的邪恶，这些人正是通过他们的结合形成婚姻。事实上，可耻情欲之恶可以在没有婚姻的情况下存在，而婚姻本来也可以没有它。然而，婚姻不能没有它而存在，尽管它可以在没有婚姻的情况下存在；这属于身体的状况（不是那生命的，而是）这死亡的。当然，那引起羞耻的肉性情欲在婚姻状态之外也存在，每当它催促人去犯奸淫、淫乱和不洁，这些完全敌对婚姻的纯洁；或者，当它没有犯这些事，因为人不允许或同意它们发生，但它仍然升起、发动、造成困扰，并且（尤其在梦中）造成它真实工作的相似物，达到它自身情感的目的。那么，这是一个甚至在婚姻状态中也不是婚姻之恶的邪恶；但它在死亡之身中已准备好了这个机制，甚至违背它自身的意愿，而这个机制对于完成它确实愿意的事是必不可少的。因此，所提及的邪恶并非来自婚姻本身的制度（它是被祝福的），而完全是来自这种情况：罪藉着一个人进入了世界，死亡藉着罪；于是死亡归于众人，因为众人都犯了罪。\BibleRef{罗 5:12}\par

% structure: p0093 source=h2 target=subsection reason=relative-heading-rank; base=h2

\subsection{第43章——福音中不能结出坏果实的良善之树并非指婚姻}

那么，他说\Quote{凭果实可认出树来}（因为我们在福音中读到主这样说过）是什么意思？难道主是用这些话谈论这个问题，而不是谈论人的两种意志，善的和恶的，称其中一个为好树，另一个为坏树，因为善行出于善意，恶行出于恶意——反之则不可能，即善行出于恶意，恶行出于善意？然而，如果我们按照他所提到的福音比喻，假设婚姻是好树，那么，我们当然必须另一方面假设淫乱是坏树。因此，如果一个人被称为婚姻的果实，即好树的好果实，那么无疑，一个人绝不可能由淫乱而生。\BibleQuote{因为坏树不能结出好果实。}\BibleRef{玛 7:18}再者，如果他说必须假设占据树的位置的不是奸淫，而是人从中出生的人性，那么这样，婚姻甚至不能代表树，只有人从中出生的人性才能代表树。因此，他从福音中取来的比喻对阐明这个问题毫无用处，因为婚姻不是那在自然生育中传递并在新生中得到补赎的罪的原因；而是第一个人的自愿过犯是原罪的原因。\Quote{您重复，}他说，\Quote{您的主张，\InnerQuote{正如罪，无论从何来源传给婴儿，是魔鬼的作为，同样，人，无论他如何出生，是天主的作为。}}是的，我说过这话，而且是千真万确的；如果此人不是佩拉纠派，而是公教徒，他在公教信仰中也会同样宣认。\par

% structure: p0095 source=h2 target=subsection reason=relative-heading-rank; base=h2

\subsection{第44章[XXVII.]——佩拉纠派争论说，如果罪由出生而来，那么所有已婚者都该受谴责}

那么，当他如此询问我们：\Quote{罪在婴儿身上因何而被找到？是因意志，还是因婚姻，还是因父母？}他的目的是什么呢？他确实说得好像他对所有这些问题都有答案，并且好像通过将一切与罪一同清除，他就能使婴儿身上没有任何可被找到的罪的来源。我恳请诸位注意他原话中的措辞：\Quote{罪在婴儿身上因何而被发现？}他说，\Quote{是因意志吗？但婴儿从未有过意志。是因婚姻吗？但这属于父母的工作，而你们先前曾宣称他们在这一行为中没有犯罪；尽管从你们后来的话中看出，你们并未真正作出这一让步。因此，婚姻必须受到谴责，因为它提供了恶的原因。然而婚姻仅仅指示了行动者的工作。因此，父母，通过他们的结合为罪提供了机会，理当受到谴责。那么，如果我们要遵循你们的意见，已婚者就被交付永罚，这就不再存疑，因为通过他们的手段，魔鬼得以在人身上实行统治。那么你们刚才所说的\InnerQuote{人是由天主所造的}又成了什么呢？因为如果通过他们的出生，恶在人身上发生，并且通过恶，魔鬼在人身上拥有权力，那么事实上你们就宣称魔鬼是人类的作者，因为人出自于出生。然而，如果你们相信人是由天主造的，并且夫妇是无辜的，那么看，你们的立场是多么不可能，即原罪是从他们而来的。}\par

% structure: p0097 source=h2 target=subsection reason=relative-heading-rank; base=h2

\subsection{第四十五章 —— 对此论点的答复：宗徒说我们都在一人内犯了罪。}

现在，宗徒给了他对所有这些问题的一个答案。他既不指责婴儿的意志（因为婴儿尚未成熟到足以犯罪），也不指责婚姻（因为婚姻就其本身而言，不仅有天主的设立，而且有天主的祝福），也不指责父母（就他们作为父母而言，他们正当地合法地结合以生育子女）；而是说：\BibleQuote{藉着一个人，罪进入了世界，藉着罪，死亡；于是死亡遍及了众人，因为众人都在他内犯了罪。}\BibleRef{罗 5:12} 如果这些人能以公教的心怀和耳朵接受这句话，他们就不会对基督的恩宠与信德抱有叛逆的情绪，也不会徒然试图将宗徒这些非常清晰显明的言辞曲解为他们的个别的和异端的含义。他们断言这段话的含义是：亚当是第一个犯罪的人，后来任何人若愿意犯罪，都在他身上找到了犯罪的榜样；因此，罪并不是从这人藉着出生传给众人，而是藉着仿效这人。然而确然的是，如果宗徒在这里是指仿效，他就会说罪不是藉着一个人，而是藉着魔鬼进入世界并传给了众人。因为关于魔鬼，经上记载：\BibleQuote{属于他那一方的人都仿效他。}\BibleRef{智 2:24} 他使用了\Quote{藉着一个人}这个说法——人类的世代当然是由他开始的——目的是向我们表明，原罪是藉着世代传给众人的。\par

% structure: p0099 source=h2 target=subsection reason=relative-heading-rank; base=h2

\subsection{第四十六章 —— 死亡的统治是什么；未来亚当的预象；众人如何因基督而成义。}

宗徒接着说的话又有什么含义呢？因为他在说了上述的话之后，接着又说：\BibleQuote{因为直到法律以前，罪恶已在世界上。} \BibleRef{罗 5:13} 这无异于说，连法律也不能除去罪恶。\BibleQuote{但是罪恶}，他接着说，\BibleQuote{在没有法律的时候，就不被算为罪。} \BibleRef{罗 5:13} 它当时存在，却不被归罪，因为没有法律来定它为罪。的确，基于同样的原则，他在另一处说：\BibleQuote{法律使人认识罪。} \BibleRef{罗 3:20} \BibleQuote{然而，}他说，\BibleQuote{从\Person{亚当}到\Person{梅瑟}，死亡作了王；} \BibleRef{罗 5:14} 也就是说，正如他先前所表达的，\BibleQuote{直到法律以前。} 这并不是说\Person{梅瑟}以后没有罪，而是因为连\Person{梅瑟}所颁布的法律也不能解除死亡的权势，当然，死亡只能借着罪作王。它的统治甚至使有死的人陷入那永远长存的第二次死亡。\BibleQuote{死亡作了王，}但是临到谁呢？\BibleQuote{甚至临到那些没有像\Person{亚当}的过犯那样犯罪的人，他是那要来者的预像。} \BibleRef{罗 5:14} 那要来者若不是\Person{基督}，还能是谁？又是什么样的预像，如果不是相反的方式？他在别处简短地表达说：\BibleQuote{如同在\Person{亚当}内众人都死了，照样在\Person{基督}内众人都要复活。} \BibleRef{格前 15:22} 一种情形是在一人内，同样另一种情形是在另一人内；这就是预像。但这预像并非在各方面都一致；因此宗徒顺着同样的思路接着说：\BibleQuote{但恩赐不如同过犯；如果因一人的过犯，众人都死了，那么天主的恩宠和那出于恩宠的恩赐，因一人\Person{耶稣基督}，就更丰盛地临到众人。} \BibleRef{罗 5:15} 但为什么说\Quote{更丰盛地}呢？岂不是因为所有借着\Person{基督}而得救的人，因\Person{亚当}的缘故经历暂时的死亡，却因\Person{基督}自己而获得永生？\BibleQuote{而且恩赐不像由一人犯罪的结果，}他说，\BibleQuote{因为审判是由一人导致定罪，但恩赐是由许多过犯导致成义。} \BibleRef{罗 5:15} \Quote{由一人}是指什么？岂非过犯？因为随后说：\BibleQuote{恩赐是由许多过犯而来。} 让这些反对者告诉我们，怎么可能\BibleQuote{由一次过犯导致定罪}，除非那已经传给所有人的一个原罪就足以导致定罪？而恩赐从许多过犯中拯救人，直到成义，因为它不但取消那源于原罪的一次过犯，也取消每个人因自己意志的冲动而添加的一切其他过犯。\BibleQuote{如果因一人的过犯，死亡就借着那人作了王，那么那些领受丰盛恩宠和义德的人，更要借着一人\Person{耶稣基督}在生命中作王。如此，因一次的过犯，众人都被定罪；照样，因一次的义行，众人也都被成义而获得生命。} \BibleRef{罗 5:17} 此后，让他们继续坚持他们虚妄的想象，主张一人不是借着遗传传递罪，而只是立了犯罪的榜样。那么，怎么能够因一人的过犯，审判就临到众人以致定罪，而不是因各人自己的许多罪呢？除非即使只有那一个罪，无需添加任何别的罪，就足以导致定罪——事实上，它确实导致所有从\Person{亚当}出生而没有在\Person{基督}内重生就死去的婴孩被定罪。那么，当他拒绝听从宗徒时，为什么还要我们回答他的问题：\Quote{罪在婴孩身上是如何发现的？——是通过意志，还是通过婚姻，还是通过他的父母？} 让他安静地聆听，听罪在婴孩身上是如何发现的。\BibleQuote{因一人的过犯，}宗徒说，\BibleQuote{众人都被定罪。} 他还说，众人都因\Person{亚当}而被定罪，众人都因\Person{基督}而成义：当然，这不是说\Person{基督}使所有在\Person{亚当}内死去的人都获得生命；而是因为他用了\Quote{众人}和\Quote{众人}这两个词，因为如同没有\Person{亚当}就没有人进入死亡，同样，没有\Person{基督}就没有人进入生命。正如我们说到一个教师，当他在城里唯一时，这个人教他们的全部学问；不是因为所有居民都上课，而是因为凡是求学的人，没有不是由他教的。事实上，宗徒后来用\Emph{许多人}来指称他先前描述的\Emph{众人}，用这两个不同的词指同一群人。\BibleQuote{因为，}他说，\BibleQuote{正如因\Person{亚当}一人的悖逆，许多人成为罪人，照样因\Person{基督}一人的服从，许多人将成为义人。} \BibleRef{罗 5:19}\par

% structure: p0101 source=h2 target=subsection reason=relative-heading-rank; base=h2

\subsection{第四十七章——圣经一再教导我们，众人在一人内都犯了罪}

他仍然追问他的问题：\Quote{罪在婴孩身上是如何发现的？} 他可以在受默感的书卷中找到答案：\BibleQuote{罪是因\Person{亚当}一人进入了世界，死又是因罪而来，于是死就临到众人，因为众人都犯了罪。} \BibleQuote{因\Person{亚当}一人的过犯，多人死了。} \BibleQuote{审判是由一人导致定罪。} \BibleQuote{因\Person{亚当}一人的过犯，死亡借着那人作了王。} \BibleQuote{因\Person{亚当}一人的过犯，审判临到众人以致定罪。} \BibleQuote{因\Person{亚当}一人的悖逆，多人成为罪人。} \BibleRef{罗 5:12} 看，\Quote{罪在婴孩身上是如何发现的。} 现在让他相信原罪，让他容许婴孩到\Person{基督}跟前，使他们得救。[XXVIII.] 他这段话说的是什么意思：\Quote{出生的人没有犯罪；生他的人没有犯罪；创造他的那位没有犯罪。在这些无辜的堡垒中，你假装罪是从哪些缺口进入的？} 为什么当有一扇敞开的门时，他却寻找隐藏的缝隙？\BibleQuote{因一人，}宗徒说；\BibleQuote{因一人的过犯，}宗徒说；\BibleQuote{因一人的悖逆，}宗徒说。他还想要什么？他要求什么更清楚的话？他期望什么被更强烈地重复？\par

% structure: p0103 source=h2 target=subsection reason=relative-heading-rank; base=h2

\subsection{第48章。原罪源于亚当的败坏意志。败坏意志从何而生。}

\Quote{如果，}他说，\Quote{罪出于意志，那败坏意志是犯罪的缘由；若出于本性，则本性是邪恶的。}我随即回答，罪确实出于意志。或许他想知道，原罪是否也是这样？我回答，当然是原罪。因为它也是从第一人的意志生出的；这样它既存在于他内，又传递给了所有人。至于他随后提出的：\Quote{若出于本性，则本性是邪恶的}，我请他若能够，如此回答：既然所有恶行都像坏树的果子一样从败坏意志生出，那就请他说出败坏意志本身——那结出恶果的坏树——究竟源于何处。若源于天使，天使是什么？无非是天主的美善作品。若源于人，人是什么？不也是天主的美善作品？事实上，败坏意志从天使的内在生出，从人的内在生出，在这邪恶生出之前，他们二者是什么？不都是天主的美善作品，具有良善可赞的本性吗？看，邪恶从美善生出；的确，它没有别的来源，只能来自美善。我称那意志为恶，它之前没有任何邪恶；当然也没有任何恶行，因为恶行只源于恶意志，如同从坏树而来。然而，恶意志从美善生出，这并非因为那美善是由美善的天主所造，而是因为它是由无中生有——不是从天主而生。那么，他的论证\Quote{若本性是天主的工程，那么魔鬼的工程就绝不能渗透天主的工程}还成立吗？我问，魔鬼的工程岂不首先在天主的工程中生出，即在那变为魔鬼的天使中？那么，若先前根本不存在恶的恶能从无中生有于天主的工程中，为什么已经存在于某处的恶不能渗透天主的工程？尤其当宗徒在经文中使用这一表述时说：\BibleQuote{于是死亡就传遍了众人}\BibleRef{罗 5:12}难道众人不是天主的工程吗？因此，罪恶传给了众人——换言之，魔鬼的工程渗透了天主的工程；或者换种说法，天主的工程所做的工程渗透了天主的工程。这就是为什么唯独天主拥有不变而全能的良善：在任何恶存在之前，祂创造了万物都是善的；并且从发生在祂所创造的善物中的一切恶中，祂使万物都受益。\par

% structure: p0105 source=h2 target=subsection reason=relative-heading-rank; base=h2

\subsection{第49章[XXIX]。在婴孩，本性来自天主，而本性的败坏来自魔鬼。}

\Quote{在一个人身上，意图被责备，源头被称扬，这是合理的；因为必须有两样事物才能容纳对立：然而在婴孩身上，只有一样事物，即本性；因为在他那里没有意志。现在这一样事物，}他说，\Quote{可归于天主或魔鬼。若本性，}他接着说道，\Quote{是来自天主的，就不可能有原恶在其中。若来自魔鬼，则没有任何根据使人能够被辩护为天主的工程。因此，持守原罪者完全是个摩尼教徒。}愿他听到与之相反的道理。在一个人身上，意志应受责备，本性应受称赞；因为要有两样事物来应用对立。而且，即使在婴孩身上，也不仅是一样事物，即由善天主所造的人的本性；他还有那败坏，正如宗徒明智所说，而非像佩拉纠、塞勒斯提乌斯或他们的任何门徒的愚昧所表述的那样，这败坏通过一人传给了众人。在这两样事物中，我们所说存在于婴孩中的，一样归于天主，另一样归于魔鬼。然而，由于（因这两样之一，即败坏）两者都处于魔鬼的权下，这确实没有产生任何矛盾；因为这并非出于魔鬼本身的权能，而是出于天主的权能。事实上，败坏被置于败坏之下，本性被置于本性之下，因为这两样甚至在魔鬼内也存在；因此，那些蒙爱和蒙拣选的人\BibleQuote{被从黑暗的权势中解救出来}\BibleRef{哥 1:13}——他们本是公义地被置于那权势下的——这清楚地显示，由善天主赐予成义者和善人的恩赐是何等伟大，祂从恶中引出善来。\par

% structure: p0107 source=h2 target=subsection reason=relative-heading-rank; base=h2

\subsection{第50章。邪恶的兴起与起源。婴孩的驱魔与嘘气，一个古老的基督徒礼仪。}

至于那段他自认为以敬虔笔触写下的经文，仿佛说：\Quote{如果本性出自天主，它里面就不可能有原罪}，另一个人岂不是甚至会在他眼里给出更敬虔的转折，这样说：\Quote{如果本性出自天主，它里面就不可能生出任何罪？}然而这并不真实。摩尼教徒确实主张这一点，他们企图将各种各样的邪恶浸入天主本身的本质，而不是祂从无中创造出的受造物。因为邪恶只从善中产生——但不是那至善且不变的善，即天主的本质，而是那由天主的智慧从无中创造出的善。因此，这就是人被称为天主工程的理由；因为若非天主的作为，人就不会成为人。但邪恶不会存在于婴孩身上，除非邪恶是由第一个人的意志所犯，并且原罪源自这样败坏的本性。因此，他所说的\Quote{主张原罪者完全是一个摩尼教徒}并不正确；相反，不相信原罪者完全是一个白拉奇主义者。因为并非从摩尼的毒害学说开始泛滥之时起，天主的教会中才首次以吹气驱魔仪式驱除准备受洗的婴孩——这礼仪表明他们若不先被解放脱离黑暗的权势，就不能被迁入基督的王国；\BibleRef{哥 1:13} 也并非在摩尼的书中我们读到\BibleQuote{人子来，是为寻找并拯救那失落的}，\BibleRef{路 19:10} 或\BibleQuote{罪藉着一人进入了世界}，\BibleRef{罗 5:12} 连同我们以上引用的其他类似经文；或者天主\BibleQuote{追讨父亲的罪于子女；}\BibleRef{出 20:5} 或\WorkTitle{圣咏}中写道：\BibleQuote{我是在罪孽中成形的，我母亲在罪恶中怀了我；}又或者，\BibleQuote{人成了虚妄：他的日子如影飞逝；}又或者，\BibleQuote{看，祢使我的日子已老，我在祢面前如同虚无；诚然，一切活着的人都完全是虚妄；}或宗徒说：\BibleQuote{受造之物被屈服于虚妄之下；}\BibleRef{罗 8:20} 或\WorkTitle{训道篇}中写道：\BibleQuote{虚而又虚，万事皆虚：人在太阳下劳作，有什么益处？}\BibleRef{训 1:2} 以及\WorkTitle{德训篇}中：\BibleQuote{从亚当子孙出离母腹之日，直到他们回归众生之母的日子，沉重的轭加在他们身上；}\BibleRef{德 40:1} 或宗徒又写着：\BibleQuote{在亚当内众人都死了；}\BibleRef{格前 15:22} 或圣约伯谈及自己的罪时说：\BibleQuote{那由妇人所生的人，生命短促，且充满怒气：他如花草凋谢，如影逝去，绝不停留。祢岂不也察看他，使他到祢面前受审判？谁能从污秽中洁净？没有一人，即使他在地上只活一日。}\BibleRef{约 14:1} 当他这里说到\Emph{污秽}时，仅仅是阅读经文就足以表明他意指\Emph{罪}。从话语中清楚看到他在说什么。相同的措辞和意义出现在\WorkTitle{匝加利亚先知书}中，在\BibleQuote{污秽的衣服}从大司祭身上被脱去，并对他说：\BibleQuote{我已除去你的罪过。}的地方。\BibleRef{匝 3:4} 那么，我认为所有这一切经文，以及其他类似含义的经文，它们指明人是在罪中和咒诅下出生的，并非在摩尼教徒的阴暗角落中被读到，而是在公教真理的阳光中。\par

% structure: p0109 source=h2 target=subsection reason=relative-heading-rank; base=h2

\subsection{第51章——称教导原罪者为摩尼教徒，就是控诉\Person{安波罗削}、\Person{西彼廉}及整个教会。}

此外，我该说什么关于那些在天主教会中盛行的神圣圣经注释家呢？他们从未试图将这些见证歪曲为异义，因为他们坚定不移地立足于我们最古老和最坚固的信德，从未被错谬的新奇所动摇。若我要收集这些见证并加以引用，任务就会太过冗长，而且我或许会显得对不可偏离的权威经典不够尊重。我仅提及至福的\Person{安波罗削}，\Person{白拉奇}曾对他的信德纯正给予如此显著的见证（如我已指出的）。然而，这位\Person{安波罗削}坚持认为，婴孩除了原罪之外，不需要基督医治恩宠的任何其他东西。但就戴着荣耀冠冕的\Person{西彼廉}而言，有谁会说他曾是，或可能曾是，一个摩尼教徒呢？他是在那毒害的异端出现在罗马世界之前就已殉道的。然而，在他关于婴孩圣洗的著作中，他如此有力地维护原罪，以至于宣告，甚至在第八天之前，若有需要，婴孩应当领受圣洗，以免他的灵魂失丧；并且他希望人们理解，婴孩能更容易地获得圣洗的宽恕，因为赦免给他的不完全是自己的罪，而是别人的罪。那么，让这位作者敢于称这些人为摩尼教徒吧；让他更进一步，在这可耻的指控下，玷污教会那最古老的传统，据此婴孩，如我所说，以吹气驱魔仪式被驱除，以便在脱离黑暗的权势——即魔鬼及其天使——之后，被迁入基督的王国。至于我们自己，我们更愿意与这些人以及植根于这古老信德的基督教会联合，忍受任何程度的诅咒和辱骂，也不愿与白拉奇主义者为伍，即使被公众赞扬的谄媚所覆盖。\par

% structure: p0111 source=h2 target=subsection reason=relative-heading-rank; base=h2

\subsection{第52章【XXX】——罪乃一切可耻之私欲偏情的根源。}

\Quote{你，}他问，\Quote{重复你的肯定：\InnerQuote{如果没有先犯罪，就不会有私欲偏情；然而，即使没有人犯罪，婚姻也会存在。}？}我从未说过，\Quote{不会有私欲偏情，}因为存在一种精神的私欲偏情，它渴慕智慧。\BibleRef{智 6:21}我的话是，\Quote{不会有\Emph{可耻的}私欲偏情。}让我的话语被重新细读，甚至是他所引用的那些，以便可以清楚地看出他是如何不诚实地对待它们的。然而，让他用任何他喜欢的名称来称呼它。我所说的是，除非人先前犯了罪，否则不会存在的东西，就是使他们在乐园中遮住腰身时感到羞耻的东西，并且每个人都会承认，如果悖逆的罪没有首先发生，就不会感觉到。现在，谁要想理解他们所感受到的，就应该考虑他们遮盖的是什么。因为用无花果树叶他们给自己做了\Quote{围裙}，而不是衣服；这些围裙或短裙在希腊语中称为\OriginalTerm{围腰布}{\GreekText{περιζώματα}}。人尽皆知，这些\Emph{peri-zomata}遮盖的是什么，有些拉丁作家用\Emph{campestria}这个词来解释。谁不知道什么人穿着这种短裙，以及这样的衣着遮盖了身体的哪些部分；甚至是罗马青年在\Emph{campus}裸体操练时所遮盖的同一部分，因此该围裙就被称为\Emph{campester}。\par

% structure: p0113 source=h2 target=subsection reason=relative-heading-rank; base=h2

\subsection{第53章 [XXXI.] — 私欲偏情不必为使结果实所必需}

他说：\Quote{因此，那种可能没有私欲偏情、没有身体运动、没有性器官必要——用你自己的说法——的婚姻，被你称为值得称赞的；而按照你的判断，如今实际发生的这种婚姻却是魔鬼的发明。因此，那些在你的梦境中可能存在的制度，你故意断言它们是好的，而圣经（当它说：\InnerQuote{因此，人要离开父母，与妻子结合，二人成为一体。}\BibleRef{创 2:24}）所指的那些，你却宣称是魔鬼般的邪恶，简言之，应被称为祸害，而不是婚姻。}毫不奇怪，我的这些贝拉基主义对手试图将我的话扭曲成他们想要的任何意思，因为他们一贯对圣经也如此行事，并且不只是在晦涩的段落，而是在其见证清晰明了的地方：事实上，这是所有其他异端者都遵循的惯习。现在，谁能做出这样的断言，说婚姻可能\Quote{没有身体运动，没有性器官的必要}？因为天主创造了性别；正如经上所写：\Quote{他造了他们男女。}\BibleRef{创 1:27}但是，那些将要结合在一起，并通过结合生育子女的人，怎么可能不移动他们的身体呢？当然，如果没有身体运动，一个人就不可能与他人有身体接触。那么，我们面前的问题不是关于身体运动——没有它就不可能有性交——而是关于生殖器官的可耻运动，这种运动当然可以不存在，而富有成果的结合仍然可以存在，如果生殖器官不是服从于情欲，而是简单地服从于意志，就像身体的其他肢体一样。即使在现在，在这个\Quote{死亡的身体}中，难道不是这样吗：命令给脚、手臂、手指、嘴唇或舌头，它们就立即按照我们意志的示意而运动？再举一个更奇妙的例子：甚至储存在泌尿容器中的液体，也听从我们的命令随意流出，即使我们没有因满溢而受压；同时，容纳液体的容器，如果处于健康状态，也能毫无困难地完成我们意志分配给它们的推进、挤压和排出内容的职能。那么，如果我们身体的生殖器官是顺从的，自然运动就会以更大的轻松和安静发生，人类受孕也会实现；除了那些违反自然秩序，并因公义的报应而受到这些肢体和器官难以驾驭之惩罚的人！这种惩罚被贞洁和纯洁的人所感受，他们无疑宁愿仅仅通过自然的欲望生育子女，而不是通过淫欲；而不贞洁的人，被这种病态的激情驱使，既爱娼妓也爱妻子，因这种肉体的惩罚而受到更重的精神悔恨的刺激。\par

% structure: p0115 source=h2 target=subsection reason=relative-heading-rank; base=h2

\subsection{第54章 [XXXII.] — 婚姻在罪存在后现在如何不同}

天主禁止我们说这人诬告我们所说的话：\Quote{像现在所成立的婚姻是魔鬼的发明。}实际上，它们与天主最初设立的婚姻完全相同。因为祂为繁衍人类而设立的这项祝福，祂并没有甚至从受审判的人身上夺走，正如祂没有剥夺他们的感官和肢体一样，这些无疑都是祂的恩赐，尽管他们因已招致的惩罚而注定要死。我所说的是婚姻，关于它曾有话说（只除了基督与教会的大圣事，这是制度所预表的）：\BibleQuote{为此，人要离开父亲和母亲，依附自己的妻子，两人成为一体。}\BibleRef{创 2:24}无疑，这是罪前所说；若无人犯罪，本可没有羞耻的情欲而完成。而现在，虽然在这死亡的身体中，没有这情欲就无法完成，但那仍然存在，使得人可以依附妻子，两人成为一体。因此，当有人声称现在的婚姻是一种事，若无人犯罪则可能是另一种事时，这并非关于其本质，而是关于某种变坏的品质。正如一个人被说成不同，尽管他其实是同一个人，当他改变生活作风变好或变坏时；因为作为义人是一种状态，作为罪人则是另一种，尽管人本身确实是同一个人。同样，没有羞耻情欲的婚姻是一种，有羞耻情欲的婚姻是另一种。然而，当女人合法地与丈夫结合，符合婚姻的真正制度，且对肉体所应尽的责任保持忠诚，没有犯奸淫的罪，从而合法地生育子女时，这实际上与天主最初设立的婚姻完全相同，尽管魔鬼藉着最初的诱惑引人犯罪，给婚姻本身造成了重创，但更确切地说，是给组成婚姻的男人和女人造成了重创，因为他诱使他们违背天主的命令——这一罪在天主的审判中被罚以人自身肢体的相互反叛。在这种婚姻状态下，虽然他们对自己的赤身露体感到羞耻，但他们无论如何也不能完全失去天主所指定的婚姻的福分。\par

% structure: p0117 source=h2 target=subsection reason=relative-heading-rank; base=h2

\subsection{第五十五章 [XXXIII.]——贪欲是一种疾病；\Quote{情欲}一词在教会语境中的含义。}

然后他从已婚者转到他们所生的人。正是关于这些人，我们必须与新的异端者进行最繁难的讨论，这涉及我们的主题。被某种来自天主的内在冲动所驱使，他做出了一些有助于解开整个症结的声明。因为他渴望激起对我们更大的仇恨，因为我们曾说过，婴儿即使从合法婚姻所生，也是生在罪中，于是他发表了以下言论：\Quote{你们确实断言，那些从未出生的人或许可能是好的；然而，那些充满世界、基督为之而死的人，你们却判定是魔鬼的工作，生来混乱，从起初就有罪。因此，}他继续道，\Quote{我已表明你们所做的无非是否认天主是实际存在的人的创造者。}我要说的是，我宣认除天主外没有谁是所有人的创造者，尽管确实所有人生在罪中，除非重生必会灭亡。的确，是魔鬼的劝说在他们心中播下的罪恶腐败，成了他们生在罪中的原因；而不是构成人的被造本性。然而，羞耻的情欲若不是一种疾病，就不会只在我们愿意时激动我们的肢体。即使夫妻合法而可敬的同居，若不是其中有疾病的情况，也不会引起脸红，避开任何目光并渴望隐秘。此外，宗徒不会禁止在这疾病中拥有妻子，若其中不存在疾病。希腊文中的短语\OriginalTerm{在欲望的情感中}{\GreekText{ἐν} \GreekText{πάθει} \GreekText{ἐπιθυμίας}}，被一些人译成拉丁文\Emph{in morbo desiderii} vel \Emph{concupiscentiæ}（欲望的病或贪欲的病）；被另外的人译成\Emph{in passione concupiscentiæ}（在贪欲的情欲中）；或者在不同抄本中有其他译法：无论如何，拉丁文对应词\Emph{passio}（情欲），尤其在教会用法中，通常被理解为带有指责的术语。\par

% structure: p0119 source=h2 target=subsection reason=relative-heading-rank; base=h2

\subsection{第五十六章——白拉奇主义者承认基督甚至为婴儿而死；\Person{尤利安}自刎于自己的剑下。}

但无论他对引人羞耻的\OriginalTerm{肉欲}{concupiscentia carnis}持何种意见，我请求你们注意他关于婴儿所说的话（我们正为他们劳苦），即他们认为婴儿需要一位救主，否则他们将毫无救恩而死。我再次重复他的话：\Quote{你声称，}他对我说，\Quote{那些从未出生的人可能本来可以是好人；然而那些充满世界、并且\Emph{基督为他们而死}的人，你判定是魔鬼的作为，生在紊乱状态中，从起初就有罪。}但愿他解开这整个争论，就如他解开这个问题的结一般！因为他会假装说这段他只论及成年人吗？为何？这里的话题是关于婴儿，关于刚出生的人；并且正是为此他激起我们对他们的仇视，因为我们判定他们从起初就有罪，因为我们宣称他们有罪，因为基督为他们而死。如果他们无罪，基督为何为他们而死？完全是从他们，是的，从他们，我们将找到原因，为何他认为应当激起对我的仇视。他问：\Quote{婴儿如何有罪，基督却为他们而死？}我们回答：不，婴儿如何无罪，既然基督为他们而死？这场争论需要一位法官来裁决。让基督作法官，让他告诉我们谁的死亡受益了？\BibleQuote{这是我的血，}他说，\BibleQuote{将为许多人倾流，以赦免罪过。}\BibleRef{玛 26:28} 也让宗徒参与审判，因为即使在宗徒内，也是基督亲自说话。论及天主圣父，他高呼：\BibleQuote{他既然没有怜惜自己的儿子，反而为我们众人把他交出了！}\BibleRef{罗 8:32} 我想他描述基督如此为我们众人交出，以至于婴儿在这件事上不与我们分离。但何必在此多言，连他也已不再争论此事？因为事实是，他不仅承认基督甚至为婴儿而死，而且他还从这个承认中责备我们，因为我们说这些同一婴儿有罪，基督是为他们而死。那么，现在让那位说基督为我们众人交出的宗徒，也告诉我们基督为何为我们交出。\BibleQuote{他被交出，}他说，\BibleQuote{是为了我们的过犯，复活是为了我们的成义。}\BibleRef{罗 4:25} 因此，既然甚至这个人都既承认又宣称，既接受又反对，婴儿也包括在基督为其交出的人之中；并且基督是为了我们的罪而被交出，那么婴儿当然必须有\OriginalTerm{原罪}{peccata originalia}，基督是为他们交出的；他必须在他们内有什么要医治的，他（正如他自己肯定的）不是为健康的人，而是为有病的人需要医生；\BibleRef{玛 9:12} 他必须有理由拯救他们，因为他来到世界，正如宗徒保禄所说，\BibleQuote{为拯救罪人；}\BibleRef{弟前 1:15} 他必须在他们内有要赦免的，因为他作证说他的血是\BibleQuote{为赦免罪过而倾流的；}\BibleRef{玛 26:28} 他必须有充足的理由寻找他们，因为他来了，如他所说，\BibleQuote{为寻找及拯救迷失了的人；}\BibleRef{路 19:10} 人子必须在他们内找到要毁灭的东西，因为他来正是为此目的，正如宗徒若望所说，\BibleQuote{为毁灭魔鬼的作为。}\BibleRef{若一 3:8} 现在，对于婴儿的这救恩，那以否认他们需要受伤者所需的医药的方式来肯定他们无罪的人，必是仇敌。\par

% structure: p0121 source=h2 target=subsection reason=relative-heading-rank; base=h2

\subsection{第五十七章 [第三十四章]——第一个人的大罪。}

现在请注意他接着所说的：\Quote{如果在罪之前，天主创造了一个让人出生的源头，而魔鬼创造了一个让父母受扰乱的源头，那么无疑圣洁必须归于出生的人，而有罪归于生产的人。然而，由于这将是对婚姻最明显的谴责；我求你将这观点从教会中间除去，并且真正相信万物是藉着耶稣基督造的，没有他，什么也没有造。}\BibleRef{若 1:3} 他这里说话，好像要使我们说，在人的实体中有某种由魔鬼创造的东西。魔鬼劝诱了恶作为罪；他并未将它创造为\OriginalTerm{本性}{natura}。无疑他劝诱了本性，因为人就是本性；因此藉着他的劝诱，他败坏了它。伤害肢体的人，当然不是创造它，而是损伤它。的确，那些加在身体上的伤口造成肢体跛行或行动困难；但它们并不影响使人成义的美德：然而那称为罪的伤口，却伤害了那曾被正义地活出的生命本身。这伤口在堕落的致命时刻被魔鬼强加，远比人们所知的罪更广更深。由此发生，我们的本性在那时那地被第一个人的大罪所败坏，不仅成为罪人，而且还产生罪人；然而那使圣洁生命的美德枯萎死去的软弱本身，实际上不是本性，而是\OriginalTerm{败坏}{corruptio}；正如不良的健康状态不是身体的实体或本性，而是紊乱；事实上，非常经常，甚至总是，父母的病态以某种方式植入，并重现于他们子女的身体中。\par

% structure: p0123 source=h2 target=subsection reason=relative-heading-rank; base=h2

\subsection{第五十八章——亚当的罪从他传至每一个出生的人，甚至由重生父母所生的人；橄榄树与野橄榄树的例证。}

但这罪在乐园中使人变得更坏，其严重程度远非我们所能判断。每个人在出生时都沾染了这罪，唯有在重生者中才得赦免。这种混乱如此之大，甚至从已经重生、罪已得赦并蒙遮盖的父母那里，也会传给他们的后代，使这些孩子被定罪，除非他们通过第一次的、肉体的出生而被束缚，再通过第二次的、属神的出生而被释放。造物主以好橄榄和野橄榄的例子，为这奇妙的事实提供了一个奇妙的范例：不仅从野橄榄的种子，甚至从好橄榄的种子中，长出的也只是野橄榄。因此，即使在那些自然出生后通过恩宠重生的人身上，也存在着这肉体的贪欲，它与心神的律法交战；然而，由于它在罪赦中得以赦免，它就不再被算作他们的罪，也绝不会造成任何伤害，除非人们同意它的冲动去行不法之事。但他们的后代，不是由属神的贪欲，而是由肉体的贪欲所生，就像从好橄榄树中长出的我们族类的野橄榄一样，通过自然出生从他们那里获得罪债，以至于除非重生，否则无法从这祸患中解脱。那么，此人为何声称我们把圣洁归给那些出生的人，而把罪过归给他们的父母呢？事实反而表明，即使父母有圣洁，原罪也固着在他们的子女身上，只有当他们重生时，这原罪才被消除。\par

% structure: p0125 source=h2 target=subsection reason=relative-heading-rank; base=h2

\subsection{第五十九章—— 伯拉纠派几乎不敢将贪欲置于犯罪之前的乐园中。}

既然如此，就让他随心所欲地思考这肉体的贪欲和那情欲吧——这情欲统治着不贞者，必须被贞洁者制服，却又被贞洁者和不贞者同感羞愧；因为我清楚地看出他很喜欢它。让他毫不犹豫地赞美他羞于称呼的东西吧；让他称它为（正如他实际上所称的）肢体的活力，且不必怕贞洁耳朵的荣誉；让他称它为肢体的力量，且不必在意无耻。让他说，如果他的羞耻允许的话，若没有人犯罪，这活力本应像乐园中的花朵一样繁茂；也无需遮盖那会被如此感动、无人应感到羞耻的东西；相反，有了妻子，它就会永远被运用而从不被抑制，以免如此巨大的快乐被如此巨大的幸福所剥夺。绝不可认为如此幸福的状态会在那样一个地方缺乏它所希望的东西，或是在心灵或身体上经历任何它不喜悦之事。因此，若情欲的冲动先于人的意志，那么意志会立即跟随它。妻子——在这种幸福状态中当然绝不应缺席——会被它催促，无论是即将怀孕还是已经怀孕；要么是孩子被怀上，要么是自然可赞的快乐得到满足——宁可让所有种子都灭亡，也不要让如此美好的贪欲的欲望失望。只要确保结合的双方不进行违反本性的彼此使用，那么（带着如此谦逊的保留）让他们随心所欲地使用为这目的而创造的生殖器官，享受快乐。但如果这本身违反本性的使用或许给他们带来快乐；如果前述可赞的情欲甚至渴望这种快乐；我怀疑他们是否会因为它甜美而追求它，还是因为它卑劣而厌恶它？如果他们追求满足，那么关于尊严的所有思考还有什么意义？如果他们厌恶它，那么好幸福中的平安安宁又在哪里？但此刻或许他的羞耻会唤醒他，他会说，这幸福状态的安宁如此巨大，这状态中可能存在的秩序如此完整，以至于肉体的贪欲从未先于他们的意志：只有当他们自己希望时，它才会出现；而且只有当他们需要生育子女时，他们才产生这愿望；结果便是，绝不会有种子被白白浪费，也绝不会有拥抱不随之带来怀孕和生产；肉体将服从意志，贪欲将在顺从中与之竞争。好吧，如果他对他所想象的幸福状态说了所有这些，他至少必须相当确定他所描述的现在在人类中并不存在。即使他不肯承认情欲是一种败坏的状态，至少让他承认，在乐园中那男人和女人的不服从之后，他们肉体的贪欲本身被败坏了，以至于曾经顺从有序地激动的东西，现在以不顺从和无序的方式激动，甚至到了这种程度：它不服从于即使是贞洁的丈夫和妻子的意志，以至于它在不需要时被激动；每当有必要时，它从不跟随他们的意志，而是有时太匆忙，有时太迟缓，行使它自己的运动。这就是这贪欲的叛逆，原祖父母因自己的不服从而承受了它，并通过自然传承传递给我们。它绝不是按他们的命令，而是在完全无序的状态下激动起来，那时他们遮盖了他们的肢体，这些肢体起初值得夸耀，但那时已成了可耻的。\par

% structure: p0127 source=h2 target=subsection reason=relative-heading-rank; base=h2

\subsection{第六十章—— 不要让伯拉纠派在对婴儿的残忍辩护中恣意妄为。}

然而，如我所说，让他对自己对这种私欲偏情的看法随心所欲；让他随意宣扬它，尽其所能赞美它（而且他\Emph{极其}喜欢，正如他几段摘录所显示的），以便佩拉纠派中那些未能享受婚姻所规定的节制之限制的人，即使不能从它的用途中得到满足，至少能从它的赞美中得到满足。只要他放过婴儿，不要无用地赞美他们的境况，也不要残忍地为他们辩护。让他不要宣称他们是安全的；让他允许他们前来，不是到\Person{佩拉纠}那里获得颂扬，而是到基督那里获得救恩。因为，为了结束这本书，既然这个人的论述已经结束（这是在你寄给我的短纸上写的），我将以他的最后几句话作结：\Quote{真正相信万物都是藉耶稣基督造的，没有祂，什么也没有造。}\BibleRef{若 1:3} 让他承认耶稣对婴儿也是耶稣；既然他承认万物都是藉祂造的，因为祂是天主圣言，那么让他也承认婴儿也是藉祂而得救，因为祂是耶稣；让他，我再说，这样做，如果他愿意做一个公教基督徒。因为福音中这样记载：\Quote{他们要称祂的名字为耶稣，因为祂要把自己的百姓从他们的罪恶中拯救出来。}\BibleRef{玛 1:21}——耶稣，因为耶稣在拉丁文中是\Emph{Salvator}，即\Quote{救主}。祂确实要拯救祂的百姓；在祂的百姓中当然有婴儿。祂要从\Quote{他们的罪恶}中拯救他们；因此，在婴儿身上也有原罪，因为祂可以成为耶稣，即救主，甚至对他们也是如此。\par

\SourceRecord{Translated by Peter Holmes and Robert Ernest Wallis, and revised by Benjamin B. Warfield.；Nicene and Post-Nicene Fathers, First Series；Vol. 5.；Edited by Philip Schaff.；Buffalo, NY: Christian Literature Publishing Co.,；1887.；<http://www.newadvent.org/fathers/15072.htm>.}
